% Problems, Chapter 21: Integral Equations
% Arfken, Weber & Harris, Mathematical Methods for Physicists, 7th ed.
% Requires: amsmath, amssymb, tikz

\section*{Problems}

\Needspace{8\baselineskip}
\subsection*{21.1 Introduction}

\Needspace{5\baselineskip}
\par\bigskip
\noindent\textbf{21.1.1}\quad Starting with the ODE, integrate
twice and derive the Volterra integral equation corresponding to
\begin{enumerate}
\item[(a)] $y''(x)-y(x)=0$;\quad $y(0)=0$,\quad $y'(0)=1$.
\par\nobreak\hfill\emph{ANS.}\enspace
$y=\displaystyle\int_{0}^{x}(x-t)y(t)\,dt+x$.
\item[(b)] $y''(x)-y(x)=0$;\quad $y(0)=1$,\quad $y'(0)=-1$.
\par\nobreak\hfill\emph{ANS.}\enspace
$y=\displaystyle\int_{0}^{x}(x-t)y(t)\,dt-x+1$.
\end{enumerate}
Check your results with Eq.~(21.11).
\solution
Integrating $y''(t)=y(t)$ from $t=0$ to $t=u$ gives
\[
y'(u)=y'(0)+\int_0^u y(t)\,dt.
\]
Integrating this from $u=0$ to $u=x$ gives
\[
y(x)=y(0)+xy'(0)+\int_0^xdu\int_0^u y(t)\,dt.
\]
For $x\ge0$ the triangular integration region is
$0\le t\le u\le x$. Reverse the order:
\[
\int_0^x y(t)\left(\int_t^xdu\right)dt
=\int_0^x(x-t)y(t)\,dt.
\]
For $x<0$ the identity
$\int_0^xdu\int_0^uy(t)\,dt=\int_0^x(x-t)y(t)\,dt$ still holds: both
sides vanish at $x=0$, and both have $x$-derivative $\int_0^xy(t)\,dt$
(the left side by the fundamental theorem of calculus, the right side by
Leibniz's rule, whose endpoint term $(x-x)y(x)$ vanishes), so they agree
for every $x$.
Thus the general initial-value formula is
\[
y(x)=y(0)+xy'(0)+\int_0^x(x-t)y(t)\,dt.
\]
Inserting the two pairs of initial data yields, respectively,
\[
\boxed{y(x)=x+\int_0^x(x-t)y(t)\,dt},\qquad
\boxed{y(x)=1-x+\int_0^x(x-t)y(t)\,dt}.
\]
To check with Eq.~(21.11), write the ODE as $y''+A(x)y'+B(x)y=g(x)$
with initial point $a$: here $A=0$, $B=-1$, $g=0$, $a=0$. Hence
\[
\begin{aligned}
K(x,t)&=(t-x)[B(t)-A'(t)]-A(t)=(t-x)(-1)=x-t,\\
f(x)&=\int_0^x(x-t)g(t)\,dt+[A(0)y(0)+y'(0)](x-0)+y(0)
=y(0)+xy'(0),
\end{aligned}
\]
exactly as above.


\Needspace{5\baselineskip}
\par\bigskip
\noindent\textbf{21.1.2}\quad Starting with the given answers of
Exercise~21.1.1, differentiate and recover the original ODEs and
the boundary conditions.
\solution
Write either answer as
\[
y(x)=a+bx+\int_0^x(x-t)y(t)\,dt,
\]
where $(a,b)=(0,1)$ or $(1,-1)$. Leibniz's rule includes an
upper-endpoint term, which vanishes here:
\[
y'(x)=b+(x-x)y(x)+\int_0^x
\frac{\partial}{\partial x}\bigl[(x-t)y(t)\bigr]dt
=b+\int_0^x y(t)\,dt.
\]
Differentiate once more; by the fundamental theorem of calculus,
$\frac{d}{dx}\int_0^xy(t)\,dt=y(x)$, so
\[
\boxed{y''(x)=y(x)}.
\]
At $x=0$ both integrals have zero length, so
\[
y(0)=a,\qquad y'(0)=b.
\]
For $(a,b)=(0,1)$ this is $y(0)=0$, $y'(0)=1$, and for $(a,b)=(1,-1)$
it is $y(0)=1$, $y'(0)=-1$: precisely the two original initial
conditions. The differentiability used here follows from the integral
equation itself: for continuous $y$ its right-hand side is continuously
differentiable, hence so is $y$, and then $y'=b+\int_0^xy\,dt$ is
differentiable again.


\Needspace{5\baselineskip}
\par\bigskip
\noindent\textbf{21.1.3}\quad Given
\[
\varphi(x)=x-\int_{0}^{x}(t-x)\varphi(t)\,dt,
\]
solve this integral equation by converting it to an ODE (plus
boundary conditions) and solving the ODE (by inspection).
\solution
First account for the sign, $-(t-x)=x-t$:
\[
\varphi(x)=x+\int_0^x(x-t)\varphi(t)\,dt.
\]
Differentiating by Leibniz's rule, the upper-endpoint term
$(x-x)\varphi(x)$ is zero; a second differentiation uses the
fundamental theorem of calculus:
\[
\varphi'(x)=1+\int_0^x\varphi(t)\,dt,\qquad
\varphi''(x)=\varphi(x).
\]
At $x=0$ the integrals in the equation and in $\varphi'$ have zero
length, so $\varphi(0)=0$, $\varphi'(0)=1$. The general solution is
$Ae^x+Be^{-x}$, and therefore
\[
A+B=0,\qquad A-B=1,\qquad
A=\frac12,\quad B=-\frac12.
\]
Hence $\boxed{\varphi(x)=\sinh x}$. Directly, integrating by parts
with $u=x-t$, $dv=\sinh t\,dt$, $v=\cosh t$, $du=-dt$,
\[
\int_0^x(x-t)\sinh t\,dt
=\left[(x-t)\cosh t\right]_0^x+\int_0^x\cosh t\,dt
=-x+\sinh x,
\]
so substitution gives $x+(-x+\sinh x)=\sinh x$.


\Needspace{5\baselineskip}
\par\bigskip
\noindent\textbf{21.1.4}\quad Show that the homogeneous Volterra
equation of the second kind
\[
\psi(x)=\lambda\int_{0}^{x}K(x,t)\psi(t)\,dt
\]
has no solution (apart from the trivial solution $\psi=0$).
\erratum{Assume that $\psi$ is continuous and $K$ is continuous
on every finite triangle $0\le t\le x\le L$.}
\par\emph{Hint.}~\erratum{On a fixed triangle, set
$M=\max|K|$, $C=\max|\psi|$, and iterate the integral bound
to obtain $|\psi(x)|\le C(|\lambda|Mx)^n/n!$.}
\solution
Use the usual continuous-kernel setting on a finite triangle
$0\le t\le x\le L$. Put
\[
M=\max_{0\le t\le x\le L}|K(x,t)|,\qquad
C=\max_{0\le x\le L}|\psi(x)|.
\]
For $0\le x\le L$ the equation gives
\[
|\psi(x)|\le|\lambda|\int_0^x|K(x,t)|\,|\psi(t)|\,dt
\le|\lambda|\int_0^xMC\,dt=C|\lambda|Mx.
\]
Suppose $|\psi(t)|\le C(|\lambda|Mt)^n/n!$. Substituting this
bound into the same estimate gives
\[
|\psi(x)|\le|\lambda|M\int_0^x
C\frac{(|\lambda|Mt)^n}{n!}\,dt
=C\frac{(|\lambda|Mx)^{n+1}}{(n+1)!}.
\]
Starting with $n=0$, induction therefore proves this bound for every
integer $n\ge0$. For fixed $x$, the ratio of successive bounds is
$|\lambda|Mx/(n+1)\to0$, so the bound tends to zero. Thus
\[
\boxed{\psi(x)=0\quad(0\le x\le L)}
\]
for every finite $\lambda$. Taking arbitrary finite $L$ covers an
unbounded interval whenever the kernel is locally bounded and continuous.

The textbook's hint instead expands $\psi$ in a Maclaurin series.
If $K$ is smooth, put $k_j(x)=\partial_x^jK(x,t)\big|_{t=x}$.
Each differentiation of $\int_0^x\partial_x^jK(x,t)\psi(t)\,dt$ by
Leibniz's rule adds the endpoint term $k_j(x)\psi(x)$, so induction on
$n$ gives
\[
\psi^{(n)}(x)=\lambda\int_0^x\partial_x^nK(x,t)\psi(t)\,dt
+\lambda\sum_{j=0}^{n-1}\frac{d^{\,n-1-j}}{dx^{n-1-j}}
\bigl[k_j(x)\psi(x)\bigr].
\]
At $x=0$ the integral has zero length, and by the product rule the sum
contains only $\psi(0),\psi'(0),\ldots,\psi^{(n-1)}(0)$. Since
$\psi(0)=0$ (zero-length integral), induction on $n$ makes every
$\psi^{(n)}(0)$ zero.
\erratum{A vanishing Maclaurin series proves the conclusion only if
$\psi$ is analytic, not merely indefinitely differentiable. For example,
the smooth extension of $e^{-1/x^2}$ by zero at $x=0$ has all derivatives
zero there but is not the zero function. The bounded-kernel estimate
above supplies the missing uniqueness argument without assuming
analyticity.}


\Needspace{8\baselineskip}
\subsection*{21.2 Some Special Methods}

\Needspace{5\baselineskip}
\par\bigskip
\noindent\textbf{21.2.1}\quad The kernel of a Fredholm equation of
the second kind,
\[
\varphi(x)=f(x)+\lambda\int_{-\infty}^{\infty}K(x,t)\varphi(t)\,dt,
\]
is of the form $k(x-t)$.\footnote{This kernel and a range
$0\le x<\infty$ are the characteristics of integral equations of
the Wiener--Hopf type. Details will be found in Chapter~8 of Morse
and Feshbach (1953); see the Additional Readings.} Assuming that
the required transforms exist, show that
\[
\varphi(x)=\frac{1}{\sqrt{2\pi}}\int_{-\infty}^{\infty}
\frac{F(t)e^{-ixt}}{1-\sqrt{2\pi}\,\lambda K(t)}\,dt.
\]
$F(t)$ and $K(t)$ are the Fourier transforms of $f(x)$ and $k(x)$,
respectively.
\erratum{Assume also that the denominator has no real zeros and
that the quotient in the displayed formula admits Fourier inversion.}
\solution
Use the textbook's symmetric Fourier convention
\[
\widehat g(\omega)=\frac1{\sqrt{2\pi}}\int_{\mathbb R}
g(x)e^{i\omega x}dx,\qquad
g(x)=\frac1{\sqrt{2\pi}}\int_{\mathbb R}
\widehat g(\omega)e^{-i\omega x}d\omega.
\]
Let $\Phi=\widehat\varphi$, $F=\widehat f$, $K=\widehat k$
(the transform variable is called $\omega$ here and $t$ in the
statement), and write
$(k*\varphi)(x)=\int_{\mathbb R}k(x-t)\varphi(t)\,dt$ for the integral
in the equation (no factor $1/\sqrt{2\pi}$). To keep track of
normalization, set $u=x-t$ for fixed $t$ in its transform, so that
$dx=du$ and $e^{i\omega x}=e^{i\omega u}e^{i\omega t}$:
\[
\begin{aligned}
\widehat{k*\varphi}(\omega)
&=\frac1{\sqrt{2\pi}}\int_{\mathbb R}\!\!\int_{\mathbb R}
k(x-t)\varphi(t)e^{i\omega x}\,dt\,dx\\
&=\frac1{\sqrt{2\pi}}
\left(\int_{\mathbb R}k(u)e^{i\omega u}du\right)
\left(\int_{\mathbb R}\varphi(t)e^{i\omega t}dt\right)\\
&=\frac1{\sqrt{2\pi}}\bigl[\sqrt{2\pi}K(\omega)\bigr]
\bigl[\sqrt{2\pi}\Phi(\omega)\bigr]
=\sqrt{2\pi}K(\omega)\Phi(\omega).
\end{aligned}
\]
Transforming each term of the integral equation therefore gives
$\Phi=F+\sqrt{2\pi}\lambda K\Phi$, that is,
\[
[1-\sqrt{2\pi}\lambda K(\omega)]\Phi(\omega)=F(\omega).
\]
Divide by the bracket and apply the inverse transform to obtain
\[
\boxed{\varphi(x)=\frac1{\sqrt{2\pi}}\int_{\mathbb R}
\frac{F(\omega)e^{-i\omega x}}
{1-\sqrt{2\pi}\lambda K(\omega)}\,d\omega}.
\]
\erratum{Existence of the transforms of $f$ and $k$ alone does not justify
division: the quotient must belong to the intended inverse-transform
class. Zeros of the denominator require compatibility with $F$ and may
permit homogeneous components supported on those zeros in a
distributional setting. The displayed formula applies when these
obstructions are absent.}
For an example of such a homogeneous component, suppose
$\sqrt{2\pi}\lambda K(\omega_0)=1$ at a real $\omega_0$. Then
$e^{-i\omega_0x}$ solves the homogeneous equation: with $u=x-t$,
\[
\lambda\int_{\mathbb R}k(x-t)e^{-i\omega_0t}dt
=\lambda e^{-i\omega_0x}\int_{\mathbb R}k(u)e^{i\omega_0u}du
=\sqrt{2\pi}\lambda K(\omega_0)e^{-i\omega_0x}=e^{-i\omega_0x}.
\]
Its transform $\sqrt{2\pi}\,\delta(\omega-\omega_0)$ is supported at
that zero, and it may be added to any solution $\varphi$.


\Needspace{5\baselineskip}
\par\bigskip
\noindent\textbf{21.2.2}
\begin{enumerate}
\item[(a)] The kernel of a Volterra equation of the first kind,
\[
f(x)=\int_{0}^{x}K(x,t)\varphi(t)\,dt,
\]
has the form $k(x-t)$. Assuming that the required transforms
exist, show that
\[
\varphi(x)=\frac{1}{2\pi i}\int_{\gamma-i\infty}^{\gamma+i\infty}
\frac{F(s)}{K(s)}\,e^{xs}\,ds,
\]
where $F(s)$ and $K(s)$ are, respectively, the Laplace
transforms of $f(x)$ and $k(x)$.
\item[(b)] In terms of the notation of part (a), show that the
Volterra equation of the second kind,
\[
\varphi(x)=f(x)+\lambda\int_{0}^{x}K(x,t)\varphi(t)\,dt,
\]
has solution
\[
\varphi(x)=\frac{1}{2\pi i}\int_{\gamma-i\infty}^{\gamma+i\infty}
\frac{F(s)}{1-\lambda K(s)}\,e^{xs}\,ds.
\]
\end{enumerate}
\erratum{Assume each quotient is a Laplace transform for which the
Bromwich inversion is valid. Choose $\gamma$ in its right half-plane
of analyticity, to the right of its singularities.}
\solution
For $x\ge0$ use
\[
F(s)=\int_0^\infty e^{-sx}f(x)dx,\quad
K(s)=\int_0^\infty e^{-sx}k(x)dx,\quad
\Phi(s)=\int_0^\infty e^{-sx}\varphi(x)dx.
\]
In a right half-plane where the integrals can be interchanged,
reverse the order on the region $0\le t\le x<\infty$ (that is,
$t\ge0$, $x\ge t$) and then put $u=x-t$, $e^{-sx}=e^{-su}e^{-st}$:
\[
\begin{aligned}
\int_0^\infty e^{-sx}\int_0^xk(x-t)\varphi(t)dt\,dx
&=\int_0^\infty\varphi(t)\left[\int_t^\infty
e^{-sx}k(x-t)\,dx\right]dt\\
&=\int_0^\infty\varphi(t)e^{-st}\left[\int_0^\infty
e^{-su}k(u)\,du\right]dt
=K(s)\Phi(s).
\end{aligned}
\]
(a) Transforming the first-kind equation gives $F=K\Phi$, hence
\[
\Phi(s)=\frac{F(s)}{K(s)},\qquad
\boxed{\varphi(x)=\frac1{2\pi i}
\int_{\gamma-i\infty}^{\gamma+i\infty}
\frac{F(s)}{K(s)}e^{sx}ds}.
\]
(b) Transforming the second-kind equation gives
$\Phi=F+\lambda K\Phi$, i.e.\ $(1-\lambda K)\Phi=F$, so
\[
\Phi(s)=\frac{F(s)}{1-\lambda K(s)},\qquad
\boxed{\varphi(x)=\frac1{2\pi i}
\int_{\gamma-i\infty}^{\gamma+i\infty}
\frac{F(s)}{1-\lambda K(s)}e^{sx}ds}.
\]
\erratum{In each case the quotient must be an admissible Laplace
transform; its singularities determine the inversion line
$\Re s=\gamma$. A zero of the divisor ($K(s)$ in (a), $1-\lambda K(s)$
in (b)) that is not cancelled by a zero of $F(s)$ is such a singularity,
since the quotient is unbounded near it. When classical inversion
hypotheses hold, choose $\gamma$ to the right of all quotient
singularities and use the ordinary Bromwich inverse.}


\Needspace{5\baselineskip}
\par\bigskip
\noindent\textbf{21.2.3}\quad Using the Laplace transform solution
(Exercise~21.2.2), solve
\begin{enumerate}
\item[(a)] $\varphi(x)=x+\displaystyle\int_{0}^{x}
(t-x)\varphi(t)\,dt$.
\par\nobreak\hfill\emph{ANS.}\enspace$\varphi(x)=\sin x$.
\item[(b)] $\varphi(x)=x-\displaystyle\int_{0}^{x}
(t-x)\varphi(t)\,dt$.
\par\nobreak\hfill\emph{ANS.}\enspace$\varphi(x)=\sinh x$.
\end{enumerate}
Check your results by substituting back into the original
integral equations.
\solution
Both equations have the form of Exercise~21.2.2(b) with $\lambda=1$
and $f(x)=x$. In (a) the kernel is $t-x=-(x-t)$, so $k(u)=-u$; in (b)
it is $-(t-x)=x-t$, so $k(u)=u$. Since $\mathcal L\{x\}=1/s^2$, the
transforms are $F=s^{-2}$ and $K=\mp s^{-2}$:
\[
\begin{array}{c|c|c}
&\text{transformed equation}&\Phi(s)\\ \hline
\text{(a)}&\Phi=s^{-2}-s^{-2}\Phi&(s^2+1)^{-1}\\
\text{(b)}&\Phi=s^{-2}+s^{-2}\Phi&(s^2-1)^{-1}
\end{array}
\]
(multiply each transformed equation by $s^2$: $(s^2\pm1)\Phi=1$).
By the standard transforms $\mathcal L\{\sin x\}=1/(s^2+1)$ and
$\mathcal L\{\sinh x\}=1/(s^2-1)$, the inverses are
\[
\boxed{\varphi_a(x)=\sin x,\qquad\varphi_b(x)=\sinh x}.
\]
For an explicit substitution check, integration by parts
(differentiate $x-t$, integrate $\sin t$ or $\sinh t$) gives
\[
\int_0^x(x-t)\sin t\,dt
=\left[-(x-t)\cos t\right]_0^x-\int_0^x\cos t\,dt
=x-\sin x,
\]
and
\[
\int_0^x(x-t)\sinh t\,dt
=\left[(x-t)\cosh t\right]_0^x+\int_0^x\cosh t\,dt
=\sinh x-x.
\]
Since $\int_0^x(t-x)\varphi\,dt=-\int_0^x(x-t)\varphi\,dt$,
equation (a) therefore reads $x-(x-\sin x)=\sin x$;
(b) reads $x+(\sinh x-x)=\sinh x$.


\Needspace{5\baselineskip}
\par\bigskip
\noindent\textbf{21.2.4}\quad Reformulate the equations of
Example~21.2.1 for integrals on the range $(0,\infty)$ using
Fourier cosine transforms.
\erratum{Take the unknown and convolution kernel to be even functions
on the full line, and retain the reflected kernel in the half-line equation.}
\solution
Take $k$ and $\varphi$ even on the full line. Their convolution
$f(x)=\int_{\mathbb R}k(x-t)\varphi(t)\,dt$ is even as well: replacing
$t$ by $-t$ in $f(-x)$ and using $k(-y)=k(y)$, $\varphi(-t)=\varphi(t)$
gives back $f(x)$. Split the full-line equation of Example~21.2.1 into
positive and negative $t$, and change the negative half-line by
$t\mapsto-t$:
\[
\int_{-\infty}^0k(x-t)\varphi(t)\,dt
=\int_0^\infty k(x+t)\varphi(-t)\,dt
=\int_0^\infty k(x+t)\varphi(t)\,dt.
\]
Hence, for $x\ge0$,
\[
\boxed{f(x)=\int_0^\infty
[k(x-t)+k(x+t)]\varphi(t)\,dt}.
\]
\erratum{A half-line cosine formulation requires these even extensions
and the reflected kernel $k(x+t)$. Simply replacing the full-line limits
by $0,\infty$ while keeping only $k(x-t)$ does not yield the same
convolution theorem.}

Define the self-inverse cosine transform
\[
G_c(\omega)=\sqrt{\frac2\pi}\int_0^\infty
g(x)\cos(\omega x)dx.
\]
For even $g$, this equals its symmetric full-line Fourier transform,
because $g(x)\sin(\omega x)$ is odd and $g(x)\cos(\omega x)$ is even:
\[
\frac1{\sqrt{2\pi}}\int_{\mathbb R}g(x)e^{i\omega x}dx
=\frac1{\sqrt{2\pi}}\int_{\mathbb R}g(x)\cos(\omega x)dx
=\frac2{\sqrt{2\pi}}\int_0^\infty g(x)\cos(\omega x)dx
=G_c(\omega).
\]
Thus the equations of the example become
\[
F_c(\omega)=\sqrt{2\pi}K_c(\omega)\Phi_c(\omega),\qquad
\Phi_c(\omega)=\frac{F_c(\omega)}{\sqrt{2\pi}K_c(\omega)}.
\]
Since $G_c$ is even in $\omega$, the same argument applied to the
inverse transform gives
\[
g(x)=\sqrt{\frac2\pi}\int_0^\infty G_c(\omega)\cos(\omega x)d\omega,
\]
so the cosine transform is its own inverse. For completeness, inversion
of the first relation gives
\[
f(x)=\sqrt{\frac2\pi}\int_0^\infty\sqrt{2\pi}K_c(\omega)\Phi_c(\omega)
\cos(\omega x)d\omega
=2\int_0^\infty K_c(\omega)\Phi_c(\omega)
\cos(\omega x)d\omega.
\]
Inverting the second, with $\sqrt{2/\pi}/\sqrt{2\pi}=1/\pi$, gives the
desired half-line solution:
\[
\boxed{\varphi(x)=\frac1\pi\int_0^\infty
\frac{F_c(\omega)}{K_c(\omega)}\cos(\omega x)d\omega}.
\]
The quotient and inversion need the same existence conditions as in
Exercise~21.2.1. If the half-line equation is instead written with one
half of the reflected sum,
$f(x)=\frac12\int_0^\infty[k(x-t)+k(x+t)]\varphi(t)\,dt$, then $f$ is
one half of the full-line convolution, so $F_c=\sqrt{\pi/2}\,K_c\Phi_c$,
and the factor $1/\pi$ in the solution becomes
$\sqrt{2/\pi}/\sqrt{\pi/2}=2/\pi$.


\Needspace{5\baselineskip}
\par\bigskip
\noindent\textbf{21.2.5}\quad Given the Fredholm integral
equation,
\[
e^{-x^{2}}=\int_{-\infty}^{\infty}e^{-(x-t)^{2}}\varphi(t)\,dt,
\]
apply the Fourier convolution technique of Example~21.2.1 to solve
for $\varphi(t)$ \erratum{as a tempered distribution}.
\solution
The symmetric Fourier transform of $e^{-x^2}$ follows by completing the
square, $-x^2+i\omega x=-(x-i\omega/2)^2-\omega^2/4$, and using
$\int_{\mathbb R}e^{-(x-i\omega/2)^2}dx=\sqrt\pi$ (the Gaussian integral
with its contour shifted parallel to the real axis):
\[
\frac1{\sqrt{2\pi}}\int_{\mathbb R}
e^{-x^2+i\omega x}dx=\frac{\sqrt\pi}{\sqrt{2\pi}}e^{-\omega^2/4}
=\frac1{\sqrt2}e^{-\omega^2/4}.
\]
Here $f(x)=e^{-x^2}$ and the kernel is $k(x-t)$ with $k(u)=e^{-u^2}$,
so $F=K=e^{-\omega^2/4}/\sqrt2$. The convolution theorem
$F=\sqrt{2\pi}K\Phi$ therefore gives
\[
\frac1{\sqrt2}e^{-\omega^2/4}
=\sqrt{2\pi}\left(\frac1{\sqrt2}e^{-\omega^2/4}\right)
\Phi(\omega).
\]
The Gaussian factor never vanishes, so
$\Phi(\omega)=1/\sqrt{2\pi}$. But this is precisely the Fourier
transform of the unit point mass at zero,
$(2\pi)^{-1/2}\int_{\mathbb R}\delta(t)e^{i\omega t}dt=1/\sqrt{2\pi}$;
equivalently, the inverse transform is
\[
\varphi(t)=\frac1{\sqrt{2\pi}}\int_{\mathbb R}
\frac1{\sqrt{2\pi}}e^{-i\omega t}d\omega
=\frac1{2\pi}\int_{\mathbb R}e^{-i\omega t}d\omega,
\]
the integral representation of the delta function:
\[
\boxed{\varphi(t)=\delta(t)}.
\]
Indeed, the original equation is checked without transforms:
\[
\int_{\mathbb R}e^{-(x-t)^2}\delta(t)\,dt=e^{-x^2}.
\]
\erratum{The answer is a distribution. There is no $L^1(\mathbb R)$
solution, since its Fourier transform would tend to zero at infinity;
nor an $L^2(\mathbb R)$ solution, since a nonzero constant transform is
not square integrable. In the tempered-distribution class the nonzero
Gaussian multiplier makes the solution unique: the transform $\Psi$ of
the difference of two solutions satisfies $e^{-\omega^2/4}\Psi=0$, so
$\langle\Psi,\chi\rangle
=\langle e^{-\omega^2/4}\Psi,\,e^{\omega^2/4}\chi\rangle=0$ for every
smooth compactly supported $\chi$ (since $e^{\omega^2/4}\chi$ is again
such a function), and therefore $\Psi=0$.}


\Needspace{5\baselineskip}
\par\bigskip
\noindent\textbf{21.2.6}\quad Solve Abel's equation,
\[
f(x)=\int_{0}^{x}\frac{\varphi(t)}{(x-t)^{\alpha}}\,dt,
\qquad0<\alpha<1,
\]
by the following method:
\begin{enumerate}
\item[(a)] Multiply both sides by $(z-x)^{\alpha-1}$ and integrate
with respect to $x$ over the range $0\le x\le z$.
\item[(b)] Reverse the order of integration and evaluate the
integral on the right-hand side (with respect to $x$) by
recognizing it as a beta function.
\end{enumerate}
\par\emph{Note.}
\[
\int_{t}^{z}\frac{dx}{(z-x)^{1-\alpha}(x-t)^{\alpha}}
=B(1-\alpha,\alpha)=\Gamma(\alpha)\Gamma(1-\alpha)
=\frac{\pi}{\sin\pi\alpha}.
\]
\solution
For locally integrable $\varphi$, the singular kernels are integrable
when $0<\alpha<1$.

(a) Multiply the equation by $(z-x)^{\alpha-1}$
and integrate over $0\le x\le z$:
\[
I(z):=\int_0^z(z-x)^{\alpha-1}f(x)dx
=\int_0^zdx\int_0^x
(z-x)^{\alpha-1}(x-t)^{-\alpha}\varphi(t)dt.
\]

(b) The integration region is the triangle $0\le t\le x\le z$.
Absolute integrability on this triangle follows from the beta integral
below applied to $|\varphi|$. Reversing the order ($t$ from $0$ to $z$,
then $x$ from $t$ to $z$),
\[
I(z)=\int_0^z\varphi(t)\left[
\int_t^z(z-x)^{\alpha-1}(x-t)^{-\alpha}dx\right]dt.
\]
Put $x=t+(z-t)u$, so $dx=(z-t)du$, $z-x=(z-t)(1-u)$, $x-t=(z-t)u$,
and $x=t,z$ correspond to $u=0,1$. The powers of $z-t$ cancel, since
$(\alpha-1)-\alpha+1=0$:
\[
\begin{aligned}
\int_t^z(z-x)^{\alpha-1}(x-t)^{-\alpha}dx
&=\int_0^1(1-u)^{\alpha-1}u^{-\alpha}du\\
&=B(1-\alpha,\alpha)
=\frac{\Gamma(1-\alpha)\Gamma(\alpha)}{\Gamma(1)}
=\frac{\pi}{\sin\pi\alpha},
\end{aligned}
\]
by the reflection formula $\Gamma(\alpha)\Gamma(1-\alpha)=\pi/\sin\pi\alpha$.
Hence
\[
I(z)=\frac{\pi}{\sin\pi\alpha}\int_0^z\varphi(t)dt.
\]
Differentiating with respect to $z$ (fundamental theorem of calculus)
and solving for $\varphi$,
\[
\boxed{\varphi(z)=\frac{\sin\pi\alpha}{\pi}
\frac{d}{dz}\int_0^z(z-x)^{\alpha-1}f(x)dx}.
\]
For merely locally integrable $\varphi$, $I(z)$ is $\pi/\sin\pi\alpha$
times the indefinite integral $\int_0^z\varphi(t)\,dt$, so by the
Lebesgue differentiation theorem $I'(z)=(\pi/\sin\pi\alpha)\varphi(z)$
for almost every $z$.

If $f$ is continuously differentiable up to zero, remove the apparent
singular endpoint before differentiating. Integrate by parts,
differentiating $f(x)$ and integrating $(z-x)^{\alpha-1}$ to
$-(z-x)^\alpha/\alpha$; the boundary term vanishes at $x=z$ because
$\alpha>0$:
\[
I(z)=\left[-\frac{f(x)(z-x)^\alpha}{\alpha}\right]_0^z
+\frac1\alpha\int_0^z(z-x)^\alpha f'(x)dx
=\frac{f(0)z^\alpha}{\alpha}
+\frac1\alpha\int_0^z(z-x)^\alpha f'(x)dx.
\]
Now differentiate: $\frac{d}{dz}\,f(0)z^\alpha/\alpha=f(0)z^{\alpha-1}$,
and by Leibniz's rule the upper-endpoint term
$(z-z)^\alpha f'(z)/\alpha$ vanishes while
$\partial_z(z-x)^\alpha/\alpha=(z-x)^{\alpha-1}$. Thus
\[
\boxed{\varphi(z)=\frac{\sin\pi\alpha}{\pi}
\left[f(0)z^{\alpha-1}+
\int_0^z(z-x)^{\alpha-1}f'(x)dx\right]}.
\]
This also proves sufficiency for such $f$: convolving either term with
$x^{-\alpha}$ and using the same beta integral yields respectively
$f(0)$ and $\int_0^x f'(y)dy=f(x)-f(0)$. Indeed, with $t=xu$,
\[
\int_0^x(x-t)^{-\alpha}t^{\alpha-1}dt=\int_0^1(1-u)^{-\alpha}u^{\alpha-1}du
=B(\alpha,1-\alpha)=\frac{\pi}{\sin\pi\alpha},
\]
and, reversing the order of integration as above,
\[
\int_0^x\frac{dt}{(x-t)^\alpha}\int_0^t(t-y)^{\alpha-1}f'(y)\,dy
=\int_0^xf'(y)\left[\int_y^x\frac{(t-y)^{\alpha-1}}{(x-t)^\alpha}dt\right]dy
=\frac{\pi}{\sin\pi\alpha}\int_0^xf'(y)\,dy,
\]
because the inner integral is the first one shifted: with
$t=y+(x-y)u$, $dt=(x-y)du$, $t-y=(x-y)u$, $x-t=(x-y)(1-u)$, the powers
of $x-y$ cancel and it again equals
$\int_0^1u^{\alpha-1}(1-u)^{-\alpha}du=\pi/\sin\pi\alpha$.
Multiplying the first result by $f(0)\sin\pi\alpha/\pi$ and the second
by $\sin\pi\alpha/\pi$, and adding, the two terms give
$f(0)+f(x)-f(0)=f(x)$.


\Needspace{5\baselineskip}
\par\bigskip
\noindent\textbf{21.2.7}\quad Given the generalized Abel equation
with $f(x)=1$,
\[
1=\int_{0}^{x}\frac{\varphi(t)}{(x-t)^{\alpha}}\,dt,
\qquad0<\alpha<1,\quad\erratum{x>0},
\]
solve for $\varphi(t)$ and verify that $\varphi(t)$ is a solution
of the given equation.
\par\nobreak\hfill\emph{ANS.}\enspace
$\varphi(t)=\dfrac{\sin\pi\alpha}{\pi}\,t^{\alpha-1}$.
\solution
Set $f=1$ in the first inversion formula of Exercise~21.2.6:
\[
\int_0^x(x-t)^{\alpha-1}dt=\left[-\frac{(x-t)^\alpha}{\alpha}\right]_0^x
=\frac{x^\alpha}{\alpha},
\]
so, since $\frac{d}{dx}(x^\alpha/\alpha)=x^{\alpha-1}$,
\[
\boxed{\varphi(x)=\frac{\sin\pi\alpha}{\pi}x^{\alpha-1}},
\qquad x>0.
\]
(The second formula of Exercise~21.2.6, with $f(0)=1$ and $f'=0$, gives
the same result.) Direct substitution, with $t=xu$, $dt=x\,du$,
$t^{\alpha-1}=x^{\alpha-1}u^{\alpha-1}$, $(x-t)^{-\alpha}=x^{-\alpha}(1-u)^{-\alpha}$,
yields
\[
\begin{aligned}
\int_0^x\frac{\varphi(t)}{(x-t)^\alpha}dt
&=\frac{\sin\pi\alpha}{\pi}
x^{\alpha-1-\alpha+1}
\int_0^1u^{\alpha-1}(1-u)^{-\alpha}du\\
&=\frac{\sin\pi\alpha}{\pi}B(\alpha,1-\alpha)
=\frac{\sin\pi\alpha}{\pi}
\Gamma(\alpha)\Gamma(1-\alpha)
=\frac{\sin\pi\alpha}{\pi}\,\frac{\pi}{\sin\pi\alpha}=1.
\end{aligned}
\]
Both endpoint singularities in the $u$ integral are integrable.
\erratum{The equation is valid for $x>0$, with right-hand limit one as
$x\downarrow0$. It cannot also be imposed as an ordinary zero-length
integral at $x=0$, which is zero; the solution itself is singular there.}


\Needspace{5\baselineskip}
\par\bigskip
\noindent\textbf{21.2.8}\quad A Fredholm equation of the first
kind has a kernel $e^{-(x-t)^{2}}$:
\[
f(x)=\int_{-\infty}^{\infty}e^{-(x-t)^{2}}\varphi(t)\,dt.
\]
\erratum{Assume a solution exists in
$L^2(\mathbb R,e^{-t^2}dt)$.} Show that the solution is
\erratum{the following series, converging in that weighted norm}:
\[
\varphi(x)=\frac{1}{\sqrt{\pi}}\sum_{n=0}^{\infty}
\frac{f^{(n)}(0)}{2^{n}n!}\,H_{n}(x),
\]
in which $H_{n}(x)$ is an $n$th-order Hermite polynomial.
\solution
Use the physicists' Hermite generating function:
\[
e^{-(x-t)^2}=e^{-t^2}e^{2tx-x^2}
=e^{-t^2}\sum_{n=0}^\infty H_n(t)\frac{x^n}{n!}.
\]
Taking the $n$th derivative with respect to $x$ at zero in the
integral equation (only the $x^n$ term of the series survives, giving
$H_n(t)$; differentiation under the integral is justified below) gives
\[
f^{(n)}(0)=\int_{\mathbb R}e^{-t^2}H_n(t)\varphi(t)dt.
\]
Now write $\varphi=\sum_n a_nH_n$. Hermite orthogonality is
\[
\int_{\mathbb R}e^{-t^2}H_n(t)H_m(t)dt
=\sqrt\pi\,2^nn!\,\delta_{nm}.
\]
Multiplying the expansion by $e^{-t^2}H_n$ and integrating therefore
gives
\[
a_n=\frac{1}{\sqrt\pi\,2^nn!}
\int_{\mathbb R}e^{-t^2}H_n(t)\varphi(t)dt
=\frac{f^{(n)}(0)}{\sqrt\pi\,2^nn!}.
\]
Consequently
\[
\boxed{\varphi(x)=\frac1{\sqrt\pi}\sum_{n=0}^\infty
\frac{f^{(n)}(0)}{2^nn!}H_n(x)}.
\]

\erratum{The textbook's unrestricted statement needs a function-space and
convergence qualification. If
$\varphi\in L^2(\mathbb R,e^{-t^2}dt)$, the formula is its convergent
weighted-$L^2$ Hermite expansion; it is not automatically a pointwise
series for arbitrary $f$.}
Here is the convergence justification. Write $z=a+ib$ and
$\|\varphi\|_w^2=\int e^{-t^2}|\varphi(t)|^2dt$.
Since $\Re(z-t)^2=(a-t)^2-b^2$,
\[
e^{t^2}\bigl|e^{-(z-t)^2}\bigr|^2=e^{t^2-2(a-t)^2+2b^2}
=e^{-t^2+4at-2\Re(z^2)},\qquad \Re(z^2)=a^2-b^2.
\]
Cauchy--Schwarz, applied to the product of $e^{-t^2/2}\varphi(t)$ and
$e^{t^2/2}e^{-(z-t)^2}$, and
$\int_{\mathbb R}e^{-t^2+4at}dt=\sqrt\pi\,e^{4a^2}$ (complete the square)
give
\begin{align*}
\left|\int_{\mathbb R}e^{-(z-t)^2}\varphi(t)dt\right|
&\le\|\varphi\|_w
\left(\int_{\mathbb R}e^{-t^2+4at-2\Re(z^2)}dt\right)^{1/2}
=\|\varphi\|_w\left(\sqrt\pi\,e^{2a^2+2b^2}\right)^{1/2}\\
&=\pi^{1/4}e^{|z|^2}\|\varphi\|_w.
\end{align*}
Derivatives of the Gaussian add only polynomial factors in $t-z$.
On a compact set of $z$, these are still dominated by an integrable
Gaussian times a polynomial after the same Cauchy--Schwarz estimate.
Differentiation under the integral is therefore valid to every order,
locally uniformly, and the convolution is entire. Parseval's identity
for the Hermite basis gives
\[
\|\varphi\|_w^2=\sum_{n=0}^\infty |a_n|^2\sqrt\pi\,2^nn!.
\]
Substituting the coefficients found above,
$|a_n|^2\sqrt\pi\,2^nn!=|f^{(n)}(0)|^2\sqrt\pi\,2^nn!/(\pi\,4^n(n!)^2)$,
gives the necessary condition
\[
\sum_{n=0}^\infty
\frac{|f^{(n)}(0)|^2}{\sqrt\pi\,2^nn!}<\infty.
\]
Conversely, if $f$ is entire and this condition holds, the orthogonal
series defines a weighted-$L^2$ function. By the moment formula above,
its Gaussian convolution has $n$th derivative
$\sqrt\pi\,2^nn!\,a_n=f^{(n)}(0)$ at zero, i.e.\ the same Taylor
coefficients as $f$; being entire, it equals $f$.

The limitation is already visible in Exercise~21.2.5:
$f(x)=e^{-x^2}$ has distributional solution $\delta$. Setting $t=0$ in
the generating function gives $e^{-x^2}=\sum_nH_n(0)x^n/n!$, so
$f^{(n)}(0)=H_n(0)$; comparing with
$e^{-x^2}=\sum_j(-1)^jx^{2j}/j!$, both vanish for odd $n$ and
$f^{(2j)}(0)=H_{2j}(0)=(-1)^j(2j)!/j!$. At $x=0$ the proposed series
therefore has the even terms
\[
\frac1{\sqrt\pi}
\frac{f^{(2j)}(0)H_{2j}(0)}{2^{2j}(2j)!}
=\frac1{\sqrt\pi}\frac{[(2j)!]^2}{(j!)^2\,4^j(2j)!}
=\frac1{\sqrt\pi}\frac{\binom{2j}{j}}{4^j}
\sim\frac1{\pi\sqrt j},
\]
where Stirling's formula gives $\binom{2j}{j}\sim4^j/\sqrt{\pi j}$.
By comparison with $\sum_j j^{-1/2}$ the positive series diverges, so
that example cannot be interpreted as an ordinary pointwise Hermite
solution.


\Needspace{5\baselineskip}
\par\bigskip
\noindent\textbf{21.2.9}\quad Solve the integral equation
\[
f(x)=\int_{-1}^{1}\frac{\varphi(t)}
{(1-2xt+x^{2})^{1/2}}\,dt,\qquad-1\le x\le1,
\]
for the unknown function $\varphi(t)$, \erratum{with $s=0,1,2,\ldots$,} if
\begin{enumerate}
\item[(a)] $f(x)=x^{2s}$,
\item[(b)] $f(x)=x^{2s+1}$.
\end{enumerate}
\par\nobreak\hfill\emph{ANS.}\enspace
(a) $\varphi(t)=\dfrac{4s+1}{2}\,P_{2s}(t)$,\quad
(b) $\varphi(t)=\dfrac{4s+3}{2}\,P_{2s+1}(t)$.
\solution
Let $N=2s$ in (a) and $N=2s+1$ in (b), with $s=0,1,\ldots$.
For $|x|<1$, the generating function converges uniformly in
$t\in[-1,1]$ (Weierstrass M-test: $|P_n(t)|\le1$ for $-1\le t\le1$, so
$|P_n(t)x^n|\le|x|^n$ there):
\[
(1-2xt+x^2)^{-1/2}=\sum_{n=0}^\infty P_n(t)x^n.
\]
Integrating termwise against $\varphi$,
\[
f(x)=\sum_{n=0}^\infty x^n\int_{-1}^1P_n(t)\varphi(t)dt,
\qquad|x|<1.
\]
Comparing the coefficients of this power series with those of
$f(x)=x^N$ (power series coefficients are unique), the equation requires
\[
\int_{-1}^1P_n(t)\varphi(t)dt=\delta_{nN}.
\]
If $\varphi=\sum_m a_mP_m$ in $L^2[-1,1]$, orthogonality,
$\int_{-1}^1P_nP_m\,dt=2\delta_{nm}/(2n+1)$, gives
\[
\frac{2a_n}{2n+1}=\delta_{nN},
\qquad a_n=\frac{2n+1}{2}\delta_{nN}.
\]
Thus, with $2N+1=4s+1$ in (a) and $4s+3$ in (b),
\[
\boxed{\varphi_a(t)=\frac{4s+1}{2}P_{2s}(t)},\qquad
\boxed{\varphi_b(t)=\frac{4s+3}{2}P_{2s+1}(t)}.
\]
Substitution gives
\[
\frac{2N+1}{2}\sum_{n=0}^\infty x^n
\int_{-1}^1P_n(t)P_N(t)dt=\frac{2N+1}{2}\,\frac{2}{2N+1}\,x^N=x^N,
\]
which checks both formulas. Completeness of the Legendre polynomials
proves uniqueness in $L^2$: the difference of two solutions has all
Legendre moments zero, hence vanishes.

The endpoint values follow by a limit of these finite-polynomial
solutions. For $1/2\le x<1$,
\[
1-2xt+x^2=(1-x)^2+2x(1-t)\ge2x(1-t)\ge1-t.
\]
Since the numerator $\varphi(t)$ is bounded, the integrand
$\varphi(t)(1-2xt+x^2)^{-1/2}$ is consequently dominated in absolute
value, uniformly in $x$, by $\max|\varphi|\,(1-t)^{-1/2}$, which is
integrable on $[-1,1]$. Dominated convergence as $x\to1^-$ gives
$f(1)=\lim_{x\to1^-}x^N=1$. Reflection ($t\to-t$, with
$\varphi(-t)=(-1)^N\varphi(t)$) gives $f(-1)=(-1)^N$, so the solutions hold
at both stated endpoints as improper integrals.


\Needspace{5\baselineskip}
\par\bigskip
\noindent\textbf{21.2.10}\quad Find the eigenvalues and
eigenfunctions of
\[
\varphi(x)=\lambda\int_{-1}^{1}(t-x)\varphi(t)\,dt.
\]
\solution
The image of the kernel is linear in $x$:
\[
\int_{-1}^1(t-x)\varphi(t)dt
=\int_{-1}^1t\varphi(t)dt-x\int_{-1}^1\varphi(t)dt.
\]
Every nonzero solution must therefore have the form $A+Bx$, with
$A=\lambda\int_{-1}^1t\varphi(t)dt$ and $B=-\lambda\int_{-1}^1\varphi(t)dt$.
The two required moments are (odd integrands drop out)
\[
\int_{-1}^1(A+Bt)dt=2A,\qquad
\int_{-1}^1t(A+Bt)dt=\frac{2B}{3}.
\]
Equating constant and linear coefficients gives
\[
A=\frac{2\lambda}{3}B,\qquad B=-2\lambda A.
\]
A nonzero solution has $A\ne0$, since $A=0$ would force
$B=-2\lambda A=0$. Eliminating $B$ yields
$A=\frac{2\lambda}{3}(-2\lambda A)=-\frac43\lambda^2A$, so
\[
1=-\frac43\lambda^2,\qquad\lambda^2=-\frac34,\qquad
\boxed{\lambda_\pm=\pm\frac{i\sqrt3}{2}}.
\]
Choose $A=1$, so $B=-2\lambda_\pm=\mp i\sqrt3$:
\[
\boxed{\varphi_+(x)=1-i\sqrt3\,x,\qquad
\varphi_-(x)=1+i\sqrt3\,x}.
\]
Their nonzero scalar multiples give the respective one-dimensional
eigenspaces. No real $\lambda$ permits a nontrivial solution.
The imaginary characteristic values are consistent with the real
antisymmetric kernel $K(x,t)=-K(t,x)$; the self-adjoint reality theorem
does not apply.


\Needspace{5\baselineskip}
\par\bigskip
\noindent\textbf{21.2.11}\quad Find the eigenvalues and
eigenfunctions of
\[
\varphi(x)=\lambda\int_{0}^{2\pi}\cos(x-t)\varphi(t)\,dt.
\]
\par\nobreak\hfill\emph{ANS.}\enspace
$\lambda_{1}=\lambda_{2}=\dfrac{1}{\pi}$,\qquad
$\varphi(x)=A\cos x+B\sin x$.
\solution
Use
\[
\cos(x-t)=\cos x\cos t+\sin x\sin t.
\]
Hence
\[
\int_0^{2\pi}\cos(x-t)\varphi(t)\,dt
=\cos x\int_0^{2\pi}\cos t\,\varphi(t)\,dt
+\sin x\int_0^{2\pi}\sin t\,\varphi(t)\,dt,
\]
a combination of $\cos x$ and $\sin x$ with constant coefficients.
Thus every nonzero solution is $\varphi(x)=A\cos x+B\sin x$.
Over $[0,2\pi]$,
\[
\int_0^{2\pi}\cos^2t\,dt=\pi,\quad
\int_0^{2\pi}\sin^2t\,dt=\pi,\quad
\int_0^{2\pi}\sin t\cos t\,dt=0.
\]
Inserting $\varphi(t)=A\cos t+B\sin t$, the two coefficient integrals
are $\pi A$ and $\pi B$, so the integral is
$\pi A\cos x+\pi B\sin x=\pi\varphi(x)$.
The equation becomes $(1-\lambda\pi)\varphi=0$, so
\[
\boxed{\lambda=\frac1\pi,\qquad
\varphi(x)=A\cos x+B\sin x,\quad(A,B)\ne(0,0)}.
\]
The characteristic value has a two-dimensional eigenspace.
A function orthogonal to both sine and cosine is annihilated by the
integral operator; it does not solve $\varphi=\lambda K\varphi$
nontrivially for any finite $\lambda$.


\Needspace{5\baselineskip}
\par\bigskip
\noindent\textbf{21.2.12}\quad Find the eigenvalues and
eigenfunctions of
\[
y(x)=\lambda\int_{-1}^{1}(x-t)^{2}y(t)\,dt.
\]
\par\emph{Hint.}~This problem may be treated by the
separable-kernel method or by a Legendre expansion.
\solution
Expanding $(x-t)^2=x^2-2xt+t^2$ and writing
$m_k=\int_{-1}^1t^ky(t)dt$,
\[
\int_{-1}^1(x-t)^2y(t)\,dt=m_2-2m_1x+m_0x^2,
\]
a quadratic polynomial in $x$. So every nonzero solution has the form
$y(x)=A+Bx+Cx^2$. Its three moments are (odd powers of $t$ integrate
to zero)
\[
m_0=2A+\frac{2C}{3},\quad
m_1=\frac{2B}{3},\quad
m_2=\frac{2A}{3}+\frac{2C}{5}.
\]
Equating the coefficients of $1,x,x^2$ in $y=\lambda(m_2-2m_1x+m_0x^2)$
gives $A=\lambda m_2$, $B=-2\lambda m_1$, $C=\lambda m_0$, that is,
\[
\begin{pmatrix}A\\B\\C\end{pmatrix}
=\lambda
\begin{pmatrix}
2/3&0&2/5\\
0&-4/3&0\\
2&0&2/3
\end{pmatrix}
\begin{pmatrix}A\\B\\C\end{pmatrix}.
\]
The odd equation is $B=-4\lambda B/3$; a solution with $B\ne0$
requires $1=-4\lambda/3$, giving
\[
\boxed{\lambda_0=-\frac34,\qquad y_0(x)=x}.
\]
For the even sector $(A,C)$, put $\mu=1/\lambda$, so that the first and
last rows read $\mu A=\frac23A+\frac25C$ and $\mu C=2A+\frac23C$.
A nonzero $(A,C)$ requires the determinant
\[
\begin{vmatrix}2/3-\mu&2/5\\2&2/3-\mu\end{vmatrix}
=(2/3-\mu)^2-\frac45=0,
\]
hence $2/3-\mu=\mp2/\sqrt5$ and $\mu_\pm=2/3\pm2/\sqrt5$. Taking $C=1$,
the last row says $2A+2/3=\mu_\pm$, so
$A=(\mu_\pm-2/3)/2=\pm1/\sqrt5$. The first row then holds as well:
$(\mu_\pm-\frac23)A=(\pm\frac2{\sqrt5})(\pm\frac1{\sqrt5})=\frac25$.
The eigenvalues are $\lambda_\pm=1/\mu_\pm$; rationalizing,
\[
\frac1{\frac23\pm\frac2{\sqrt5}}
=\frac{3\sqrt5}{2\sqrt5\pm6}
=\frac{3\sqrt5\,(2\sqrt5\mp6)}{20-36}
=\frac{30\mp18\sqrt5}{-16}
=\frac{-15\pm9\sqrt5}{8}.
\]
Therefore
\[
\boxed{\begin{aligned}
\lambda_+&=\frac{-15+9\sqrt5}{8},&
y_+(x)&=x^2+\frac1{\sqrt5},\\
\lambda_-&=\frac{-15-9\sqrt5}{8},&
y_-(x)&=x^2-\frac1{\sqrt5}.
\end{aligned}}
\]
All three eigenpairs satisfy every row of the matrix equation.
Since the image of the integral operator contains only quadratic
polynomials, the calculation exhausts all nonzero
finite-characteristic-value solutions.

In the Legendre method of the hint, $x^2=\frac13P_0+\frac23P_2(x)$
gives
\[
(x-t)^2=\frac23+\frac23\bigl[P_2(x)+P_2(t)\bigr]-2P_1(x)P_1(t).
\]
For even $y=a_0P_0+a_2P_2$, orthogonality
($\int_{-1}^1P_n^2dt=2/(2n+1)$) gives
\[
\int_{-1}^1(x-t)^2y(t)\,dt=\frac{4a_0}{3}+\frac{4a_2}{15}
+\frac{4a_0}{3}P_2(x),
\]
and comparing the coefficients of $P_2$ and $P_0$ in
$y=\lambda\int_{-1}^1(x-t)^2y(t)\,dt$,
\[
a_2=\frac{4\lambda}{3}a_0,\qquad
a_0=\lambda\left(\frac{4a_0}{3}+\frac{4a_2}{15}\right).
\]
With $a_0=1$, $y_\pm$ are proportional to $P_0+(4\lambda_\pm/3)P_2$,
and the second equation becomes
$1=\lambda\bigl(\frac43+\frac{16\lambda}{45}\bigr)$, i.e.\
$16\lambda^2+60\lambda-45=0$, whose roots
$(-60\pm36\sqrt5)/32$ are again $\lambda_\pm$. The odd function $P_1$
gives $\int_{-1}^1(x-t)^2P_1(t)\,dt=-\frac43P_1(x)$, i.e.\ $\lambda_0$.
To compare with the monomial form: since $P_2=(3x^2-1)/2$,
$P_0+(4\lambda_\pm/3)P_2$ has constant term
$1-2\lambda_\pm/3$ and $x^2$ coefficient $2\lambda_\pm$, whose ratio is
$1/(2\lambda_\pm)-1/3=\mu_\pm/2-1/3=\pm1/\sqrt5$, as for $y_\pm$.


\Needspace{5\baselineskip}
\par\bigskip
\noindent\textbf{21.2.13}\quad Use the separable-kernel technique
to show that
\[
\psi(x)=\lambda\int_{0}^{\pi}\cos x\sin t\,\psi(t)\,dt
\]
has no solution (apart from $\psi=0$). Explain this result in
terms of separability and symmetry.
\solution
Put
\[
c=\int_0^\pi\sin t\,\psi(t)dt.
\]
Separability forces $\psi(x)=\lambda c\cos x$. Substituting this
form into the definition of $c$,
\[
c=\lambda c\int_0^\pi\sin t\cos t\,dt
=\frac{\lambda c}{2}\,[\sin^2t]_0^\pi=0.
\]
Then $\psi(x)=\lambda c\cos x$ vanishes identically: $\boxed{\psi=0}$
for every finite $\lambda$.

The cancellation also follows from reflection about $\pi/2$:
$\sin(\pi-t)=\sin t$ and $\cos(\pi-t)=-\cos t$, so the product
has zero integral. In operator language, $K$ has rank one but
\[
K^2f(x)=\cos x
\left(\int_0^\pi\sin t\cos t\,dt\right)
\left(\int_0^\pi\sin u\,f(u)du\right)=0.
\]
The operator itself is nonzero, since $K(\sin t)=(\pi/2)\cos x$.
Its kernel is not symmetric under $x\leftrightarrow t$, so a nonzero
nilpotent operator is allowed. A nonzero \erratum{real} symmetric rank-one kernel
$c\,u(x)u(t)$ would instead have nonzero operator eigenvalue
$c\int u^2$, and hence a nontrivial finite characteristic value.
(For complex $u$ the qualification matters: $u(x)=e^{2ix}$ gives
$\int_0^\pi u^2dx=0$, and the symmetric kernel $u(x)u(t)$ is again
nilpotent.)


\Needspace{5\baselineskip}
\par\bigskip
\noindent\textbf{21.2.14}\quad Given
\[
\varphi(x)=\lambda\int_{0}^{1}(1+xt)\varphi(t)\,dt,
\]
solve for the eigenvalues and the eigenfunctions by the
separable-kernel technique.
\solution
Introduce the two moments
\[
c_0=\int_0^1\varphi(t)dt,\qquad
c_1=\int_0^1t\varphi(t)dt.
\]
Since $1+xt$ splits as $1\cdot1+x\cdot t$, the equation is
$\varphi(x)=\lambda(c_0+xc_1)$. Inserting it
into those moment definitions, with $\int_0^1dt=1$,
$\int_0^1t\,dt=\frac12$, $\int_0^1t^2dt=\frac13$, gives
\[
c_0=\lambda\left(c_0+\frac{c_1}{2}\right),\qquad
c_1=\lambda\left(\frac{c_0}{2}+\frac{c_1}{3}\right).
\]
A nonzero moment vector requires
\[
\begin{aligned}
0=\begin{vmatrix}1-\lambda&-\lambda/2\\
-\lambda/2&1-\lambda/3\end{vmatrix}
&=(1-\lambda)(1-\lambda/3)-\lambda^2/4\\
&=1-\frac\lambda3-\lambda+\frac{\lambda^2}3-\frac{\lambda^2}4
=1-\frac43\lambda+\frac1{12}\lambda^2.
\end{aligned}
\]
Multiplying by $12$, $\lambda^2-16\lambda+12=0$, and the quadratic
formula, $\lambda=8\mp\sqrt{64-12}=8\mp\sqrt{52}$, gives
\[
\boxed{\lambda_\pm=8\mp2\sqrt{13}}.
\]
The constant term of a nonzero eigenfunction cannot vanish: otherwise
$\varphi=Bx$ would require $0=\lambda c_0=\lambda B/2$, forcing $B=0$.
Normalize that term to one, writing $\varphi_\pm=1+b_\pm x$.
Then $c_0=1+b/2$, and the constant equation $1=\lambda(1+b/2)$ yields
$b=2/\lambda-2$. Rationalizing,
\[
\frac2{\lambda_\pm}=\frac{2}{8\mp2\sqrt{13}}
=\frac{2(8\pm2\sqrt{13})}{64-52}=\frac{4\pm\sqrt{13}}3,
\qquad b_\pm=\frac{4\pm\sqrt{13}}3-2,
\]
hence
\[
\boxed{\varphi_\pm(x)=1+\frac{-2\pm\sqrt{13}}3x}.
\]
Numerically the pairs are
$(0.7888974491,\ 1+0.5351837585x)$ and
$(15.2111025509,\ 1-1.8685170918x)$.
The $x$ coefficient of the equation,
$b=\lambda c_1=\lambda\bigl(\frac12+\frac b3\bigr)$, holds as well. With
$b=2/\lambda-2$ its right side is
\[
\frac\lambda2+\frac\lambda3\left(\frac2\lambda-2\right)
=\frac23-\frac\lambda6,
\]
and $2/\lambda-2=2/3-\lambda/6$ becomes, after multiplication by
$6\lambda$, $12-12\lambda=4\lambda-\lambda^2$, i.e.\
$\lambda^2-16\lambda+12=0$, which $\lambda_\pm$ satisfy. Since the image
of $K$ is two-dimensional, these are all the nontrivial eigenspaces.


\Needspace{5\baselineskip}
\par\bigskip
\noindent\textbf{21.2.15}\quad Knowing the form of the solutions
of an integral equation can be a great advantage. For
\[
\varphi(x)=\lambda\int_{0}^{1}(1+xt)\varphi(t)\,dt,
\]
assume $\varphi(x)$ to have the form $1+bx$. Substitute into the
integral equation. Integrate and solve for $b$ and $\lambda$.
\solution
Substitute the proposed function before making any spectral
assumption. Since $(1+xt)(1+bt)=1+bt+xt+bxt^2$, integrating term by term
gives
\[
\int_0^1(1+xt)(1+bt)dt
=1+\frac b2+x\left(\frac12+\frac b3\right).
\]
Equating coefficients in
$1+bx=\lambda[1+b/2+x(1/2+b/3)]$ gives
\[
1=\lambda(1+b/2),\qquad
b=\lambda(1/2+b/3).
\]
The first equation rules out $b=-2$. Eliminating $\lambda$ (multiply
the second equation by $1+b/2$ and use the first),
\[
b(1+b/2)=1/2+b/3,
\quad \frac{b^2}{2}+\frac{2b}{3}-\frac12=0,
\quad 3b^2+4b-3=0.
\]
The quadratic formula therefore gives
\[
b_\pm=\frac{-4\pm\sqrt{16+36}}6
=\frac{-4\pm2\sqrt{13}}6
=\frac{-2\pm\sqrt{13}}3.
\]
Finally, $1+b_\pm/2=(4\pm\sqrt{13})/6$, and rationalizing with
$(4\pm\sqrt{13})(4\mp\sqrt{13})=16-13=3$,
\[
\lambda_\pm=\frac1{1+b_\pm/2}
=\frac6{4\pm\sqrt{13}}
=\frac{6(4\mp\sqrt{13})}{3}
=\boxed{8\mp2\sqrt{13}}.
\]
Thus $\boxed{\varphi_\pm(x)=1+b_\pm x}$, with arbitrary
nonzero overall scale. The ansatz loses no eigenfunctions:
the integral's image is linear, and an eigenfunction with zero
constant term must vanish, as the constant coefficient equation shows.


\Needspace{5\baselineskip}
\par\bigskip
\noindent\textbf{21.2.16}\quad The equation
\[
f(x)=\int_{a}^{b}K(x,t)\varphi(t)\,dt
\]
has a degenerate kernel
$K(x,t)=\displaystyle\sum_{i=1}^{n}M_{i}(x)N_{i}(t)$.
\begin{enumerate}
\item[(a)] Show that this integral equation has no solution unless
$f(x)$ can be written as
\[
f(x)=\sum_{i=1}^{n}f_{i}M_{i}(x),
\]
where the $f_{i}$ are constants.
\item[(b)] Show that to any solution $\varphi(x)$ we may add
$\psi(x)$, provided that $\psi(x)$ has \erratum{zero bilinear
moments against every $N_i(x)$}:
\[
\int_{a}^{b}N_{i}(x)\psi(x)\,dx=0\qquad\text{for all }i.
\]
\end{enumerate}
\solution
(a) The finite sum may be brought outside the integral:
\[
\int_a^bK(x,t)\varphi(t)dt
=\sum_{i=1}^n M_i(x)
\underbrace{\int_a^bN_i(t)\varphi(t)dt}_{c_i}.
\]
Each $c_i$ is independent of $x$. Hence the equation requires
\[
\boxed{f\in\operatorname{span}\{M_1,\ldots,M_n\}},
\qquad f(x)=\sum_i c_iM_i(x),
\]
i.e.\ the stated form with constants $f_i=c_i=\int_a^bN_i\varphi\,dt$.
This is a necessary condition; to obtain a solution, the moment
equations $\int_a^bN_i\varphi\,dt=f_i$ must also be compatible.

(b) If every specified moment of $\psi$ is zero, then
\[
K(\varphi+\psi)(x)
=K\varphi(x)+\sum_iM_i(x)\int_a^bN_i(t)\psi(t)dt
=f(x)+0.
\]
Thus every such $\boxed{\varphi+\psi}$ is another solution.

\erratum{The Instructor's Manual adds that $\psi$ must itself be a linear
combination of the $M_i$. That extra restriction is false for this
first-kind equation. Only the image $f$, not the unknown or its
nullspace additions, is constrained to that span.}
For a direct counterexample, take $K(x,t)=1$ on $[0,1]$,
$f=1$, $\varphi=1$, and $\psi(t)=t-\tfrac12$. Then
\[
\int_0^1\psi(t)dt=0,\qquad
\int_0^1[1+\psi(t)]dt=1,
\]
although $\psi$ is not in the span of the constant $M_1=1$.
The displayed conditions $\int_a^bN_i\psi\,dx=0$ are exactly what the
calculation uses. They are called bilinear because no complex conjugate
is taken: for complex $N_i$ they differ from Hermitian orthogonality,
$\int_a^bN_i^*\psi\,dx=0$, which would not make the sum
$\sum_iM_i\int_a^bN_i\psi\,dt$ vanish. For example, on $[0,2\pi]$,
$N_1=e^{it}$ and $\psi=e^{-it}$ give $\int_0^{2\pi}N_1^*\psi\,dt
=\int_0^{2\pi}e^{-2it}dt=0$ but $\int_0^{2\pi}N_1\psi\,dt=2\pi$.


\Needspace{5\baselineskip}
\par\bigskip
\noindent\textbf{21.2.17}\quad A Kirchhoff diffraction theory
analysis of a laser leads to the integral equation
\[
v(\mathbf{r}_{2})=\gamma\iint
K(\mathbf{r}_{1},\mathbf{r}_{2})v(\mathbf{r}_{1})\,dA.
\]
The unknown, $v(\mathbf{r}_{1})$, gives the geometric distribution
of the radiation field over one mirror surface; the range of
integration is over the surface of that mirror. For square
confocal spherical mirrors, the integral equation becomes
\[
v(x_{2},y_{2})=\frac{-i\gamma e^{ikb}}{\lambda b}
\int_{-a}^{a}\!\!\int_{-a}^{a}
e^{-(ik/b)(x_{1}x_{2}+y_{1}y_{2})}\,v(x_{1},y_{1})\,dx_{1}\,dy_{1},
\]
in which $b$ is the centerline distance between the laser mirrors.
This can be put in a somewhat simpler form by the substitutions
\[
\frac{kx_{i}^{2}}{b}=\xi_{i}^{2},\qquad
\frac{ky_{i}^{2}}{b}=\eta_{i}^{2},\qquad\text{and}\qquad
\frac{ka^{2}}{b}=\frac{2\pi a^{2}}{\lambda b}=\alpha^{2}.
\]
\begin{enumerate}
\item[(a)] Show that the variables separate and we get two
integral equations.
\item[(b)] Show that the new limits, $\pm\alpha$, may be
approximated by $\pm\infty$ \erratum{for fixed low-order modes when
$a^2\gg\lambda b$. Explain why $a\gg\lambda$ alone is insufficient.}
\item[(c)] Solve the resulting integral equations.
\end{enumerate}
\solution
(a) Choose signed scaled coordinates
\[
\xi=\sqrt{k/b}\,x,\qquad \eta=\sqrt{k/b}\,y,\qquad
\alpha=a\sqrt{k/b}.
\]
This choice retains the signs that were suppressed by the squared
relations. The Jacobian is $dx_1dy_1=(b/k)d\xi_1d\eta_1$, the exponent
becomes $(k/b)(x_1x_2+y_1y_2)=\xi_1\xi_2+\eta_1\eta_2$, the limits
$\pm a$ become $\pm\alpha$, and, since $k\lambda=2\pi$, the prefactor
becomes
\[
\frac{-i\gamma e^{ikb}}{\lambda b}\cdot\frac bk
=\frac{-i\gamma e^{ikb}}{k\lambda}
=\frac{-i\gamma e^{ikb}}{2\pi}.
\]
The equation becomes
\[
v(\xi_2,\eta_2)=\frac{-i\gamma e^{ikb}}{2\pi}
\int_{-\alpha}^{\alpha}\!\!\int_{-\alpha}^{\alpha}
e^{-i(\xi_1\xi_2+\eta_1\eta_2)}
v(\xi_1,\eta_1)d\xi_1d\eta_1.
\]
Define
\[
(\mathcal F_\alpha h)(u)=\frac1{\sqrt{2\pi}}
\int_{-\alpha}^{\alpha}e^{-ius}h(s)ds.
\]
For a separated mode $v(\xi,\eta)=X(\xi)Y(\eta)$ the exponential
factorizes, $e^{-i(\xi_1\xi_2+\eta_1\eta_2)}=e^{-i\xi_1\xi_2}e^{-i\eta_1\eta_2}$,
so the double integral is the product
$\sqrt{2\pi}(\mathcal F_\alpha X)(\xi_2)\cdot\sqrt{2\pi}(\mathcal F_\alpha Y)(\eta_2)$
and the factor $2\pi$ cancels. The equation is
\[
XY=-i\gamma e^{ikb}(\mathcal F_\alpha X)
(\mathcal F_\alpha Y).
\]
Where $X(\xi)Y(\eta)\ne0$, divide by $XY$:
\[
\frac{(\mathcal F_\alpha X)(\xi)}{X(\xi)}\cdot
\frac{(\mathcal F_\alpha Y)(\eta)}{Y(\eta)}
=\frac1{-i\gamma e^{ikb}}.
\]
The first quotient depends only on $\xi$ and the second only on $\eta$;
holding $\eta$ fixed and varying $\xi$ shows that the first is a
constant $\rho$, and likewise the second is a constant $\sigma$.
Thus the two one-dimensional equations and their consistency
condition are
\[
\boxed{\mathcal F_\alpha X=\rho X,\qquad
\mathcal F_\alpha Y=\sigma Y,\qquad
1=-i\gamma e^{ikb}\rho\sigma}.
\]
These identities also make sense at isolated zeros of $X,Y$ by
continuity, without dividing there.

(b) \erratum{The textbook's condition $a\gg\lambda$ alone does not imply
that $\alpha$ is large: the mirror spacing $b$ also enters. The needed
large-aperture condition for fixed low-order modes is
$a^2\gg\lambda b$, equivalently $\alpha^2=2\pi a^2/(\lambda b)\gg1$.}
For example, $a=10^3\lambda$ and $b=10^9\lambda$ give
$\alpha^2=2\pi\cdot10^6\lambda^2/(10^9\lambda^2)=2\pi/1000$, much less
than one despite $a\gg\lambda$.

Let $\mathcal F$ denote $\mathcal F_\alpha$ with $\alpha=\infty$, the
full Fourier transform with kernel $e^{-ius}$. The replacement by
infinite limits is an approximation for modes
concentrated well inside $(-\alpha,\alpha)$. Since $|e^{-ius}|=1$, its
pointwise error obeys
\[
|(\mathcal F-\mathcal F_\alpha)h(u)|
=\frac1{\sqrt{2\pi}}\left|\int_{|s|>\alpha}e^{-ius}h(s)ds\right|
\le\frac1{\sqrt{2\pi}}\int_{|s|>\alpha}|h(s)|ds.
\]
For a fixed Hermite--Gaussian mode the tail is a polynomial times
$e^{-s^2/2}$ and becomes exponentially small as $\alpha\to\infty$.
For the ground Gaussian $h_0(s)=\pi^{-1/4}e^{-s^2/2}$, at $u=0$,
\[
\mathcal Fh_0(0)=\frac{\pi^{-1/4}}{\sqrt{2\pi}}
\int_{\mathbb R}e^{-s^2/2}ds=\pi^{-1/4},\qquad
(\mathcal F-\mathcal F_\alpha)h_0(0)
=\frac{2\pi^{-1/4}}{\sqrt{2\pi}}\int_\alpha^\infty e^{-s^2/2}ds
=\pi^{-1/4}\operatorname{erfc}\frac{\alpha}{\sqrt2},
\]
using $\int_\alpha^\infty e^{-s^2/2}ds=\sqrt{\pi/2}\,
\operatorname{erfc}(\alpha/\sqrt2)$ (put $s=\sqrt2\,w$ and use
$\operatorname{erfc}z=\frac2{\sqrt\pi}\int_z^\infty e^{-w^2}dw$). So the
relative error at $u=0$ is exactly
$\operatorname{erfc}(\alpha/\sqrt2)$. Higher orders require the aperture
to exceed their characteristic width as well.

(c) Now solve the infinite-limit equations. The normalized
Hermite functions are
\[
h_n(u)=\frac{H_n(u)e^{-u^2/2}}
{\pi^{1/4}\sqrt{2^nn!}}.
\]
The Gaussian transform gives $\mathcal Fh_0=h_0$: completing the
square, $-s^2/2-ius=-(s+iu)^2/2-u^2/2$, and
$\int_{\mathbb R}e^{-(s+iu)^2/2}ds=\sqrt{2\pi}$, so
\[
\mathcal Fh_0(u)=\frac{\pi^{-1/4}e^{-u^2/2}}{\sqrt{2\pi}}
\int_{\mathbb R}e^{-(s+iu)^2/2}ds=\pi^{-1/4}e^{-u^2/2}=h_0(u).
\]
To obtain all higher orders, use the raising operator $R=(u-d/du)/\sqrt2$
(applied to functions of the variable $u$, or of $s$ inside $\mathcal F$). The
Hermite recurrences $H_n'=2nH_{n-1}$ and $H_{n+1}=2uH_n-2nH_{n-1}$ give
$H_{n+1}=2uH_n-H_n'$. Write $h_n=c_nH_ne^{-u^2/2}$, where
$c_n=(\pi^{1/4}\sqrt{2^nn!})^{-1}$; by the product rule
$h_n'=c_n(H_n'-uH_n)e^{-u^2/2}$, so
\[
uh_n-h_n'=c_n(2uH_n-H_n')e^{-u^2/2}=c_nH_{n+1}e^{-u^2/2}
=\frac{c_n}{c_{n+1}}\,h_{n+1}=\sqrt{2(n+1)}\,h_{n+1},
\]
since $c_n/c_{n+1}=\sqrt{2^{n+1}(n+1)!/(2^nn!)}=\sqrt{2(n+1)}$.
Dividing by $\sqrt2$ gives $R h_n=\sqrt{n+1}\,h_{n+1}$, i.e.\
$h_n=R^nh_0/\sqrt{n!}$. Differentiation under the integral
and integration by parts (boundary terms vanish for decaying $h$) give
\[
\frac{d}{du}\mathcal Fh=\frac1{\sqrt{2\pi}}\int_{\mathbb R}(-is)e^{-ius}h(s)ds
=-i\mathcal F(sh),\qquad
\mathcal F(h')=-\frac1{\sqrt{2\pi}}\int_{\mathbb R}(-iu)e^{-ius}h(s)ds,
\]
that is,
\[
\mathcal F(sh)=i\frac{d}{du}\mathcal Fh,\qquad
\mathcal F(h')=iu\,\mathcal Fh.
\]
Therefore
\[
\mathcal F(R h)
=\frac1{\sqrt2}\bigl[\mathcal F(sh)-\mathcal F(h')\bigr]
=\frac{i}{\sqrt2}\left(\frac{d}{du}-u\right)\mathcal Fh
=-i\,R\mathcal Fh,
\]
and induction,
$\mathcal Fh_{n+1}=\mathcal F(R h_n)/\sqrt{n+1}
=-i\,R(-i)^nh_n/\sqrt{n+1}=(-i)^{n+1}h_{n+1}$, yields
\[
\boxed{\mathcal Fh_n=(-i)^nh_n}.
\]
A separated two-dimensional mode $X=h_n$, $Y=h_m$ thus has
$\rho=(-i)^n$, $\sigma=(-i)^m$, and the consistency condition
$1=-i\gamma e^{ikb}(-i)^{n+m}$ gives
$\gamma=e^{-ikb}/(-i)^{n+m+1}=i^{\,n+m+1}e^{-ikb}$ (since $1/(-i)=i$):
\[
\boxed{v_{nm}(\xi,\eta)=C\,h_n(\xi)h_m(\eta),\qquad
\gamma_{nm}=i^{\,n+m+1}e^{-ikb}}.
\]
Indeed $-i\gamma_{nm}e^{ikb}(-i)^{n+m}=1$.
In physical coordinates the shape is proportional to
\[
H_n\!\left(\sqrt{k/b}\,x\right)
H_m\!\left(\sqrt{k/b}\,y\right)
e^{-k(x^2+y^2)/(2b)}.
\]
Hermite functions form a complete $L^2(\mathbb R)$ basis.
The full one-dimensional eigenspace for phase $(-i)^r$ consists of
convergent combinations of $h_n$ with $n\equiv r\pmod4$.
Likewise, two-dimensional combinations with fixed $n+m\bmod4$
share the same $\gamma$. These are the infinite-aperture solutions;
for finite $\alpha$ the truncated transform must be solved instead,
and the tail bound quantifies the stated approximation for a chosen mode.


\Needspace{8\baselineskip}
\subsection*{21.3 Neumann Series}

\Needspace{5\baselineskip}
\par\bigskip
\noindent\textbf{21.3.1}\quad Using the Neumann series, solve
\begin{enumerate}
\item[(a)] $\varphi(x)=1-2\displaystyle\int_{0}^{x}
t\varphi(t)\,dt$,
\par\nobreak\hfill\emph{ANS.}\enspace(a) $\varphi(x)=e^{-x^{2}}$.
\item[(b)] $\varphi(x)=x+\displaystyle\int_{0}^{x}
(t-x)\varphi(t)\,dt$,
\item[(c)] $\varphi(x)=x-\displaystyle\int_{0}^{x}
(t-x)\varphi(t)\,dt$.
\end{enumerate}
\solution
(a) Define $(Tg)(x)=-2\int_0^x t\,g(t)dt$, so the equation is
$\varphi=1+T\varphi$. Iterating $\varphi_{k+1}=1+T\varphi_k$ from
$\varphi_0=1$ gives the Neumann series $\sum_{n\ge0}T^n1$.
We find its individual terms by induction:
if $T^n1=(-1)^nx^{2n}/n!$, then
\[
T^{n+1}1=-2\int_0^x t\,\frac{(-1)^nt^{2n}}{n!}dt
=-\frac{2(-1)^n}{n!}\,\frac{x^{2n+2}}{2n+2}
=\frac{(-1)^{n+1}x^{2n+2}}{(n+1)!},
\]
since $(2n+2)\,n!=2\,(n+1)!$.
The initial case is $T^01=1$, so, by the exponential series with
argument $-x^2$,
\[
\boxed{\varphi_a(x)=\sum_{n=0}^\infty
\frac{(-1)^nx^{2n}}{n!}=e^{-x^2}}.
\]
Substitution gives, since $-2te^{-t^2}=\frac{d}{dt}e^{-t^2}$,
$1-2\int_0^xt e^{-t^2}dt=1+\bigl[e^{-t^2}\bigr]_0^x
=1+(e^{-x^2}-1)=e^{-x^2}$.

(b) Put $(Jg)(x)=\int_0^x(x-t)g(t)dt$. Since $t-x=-(x-t)$, the
equation is $\varphi=x-J\varphi$. The elementary monomial formula is
\[
\begin{aligned}
J(t^p)(x)&=x\int_0^xt^p\,dt-\int_0^xt^{p+1}dt
=x\frac{x^{p+1}}{p+1}-\frac{x^{p+2}}{p+2}\\
&=\frac{x^{p+2}}{(p+1)(p+2)},\qquad p>-1.
\end{aligned}
\]
With $p=2n+1$ it maps $x^{2n+1}/(2n+1)!$ to
$x^{2n+3}/[(2n+1)!\,(2n+2)(2n+3)]=x^{2n+3}/(2n+3)!$. Starting with
$J^0x=x$, induction therefore gives $J^nx=x^{2n+1}/(2n+1)!$.
The Neumann series of $\varphi=x-J\varphi$ is
$\sum_n(-J)^nx=\sum_n(-1)^nJ^nx$. Hence
\[
\boxed{\varphi_b(x)=\sum_{n=0}^\infty
\frac{(-1)^nx^{2n+1}}{(2n+1)!}=\sin x}.
\]
(c) Here $-(t-x)=x-t$, so the equation is $\varphi=x+J\varphi$; its
Neumann series $\sum_nJ^nx$ has the same terms with positive signs:
\[
\boxed{\varphi_c(x)=\sum_{n=0}^\infty
\frac{x^{2n+1}}{(2n+1)!}=\sinh x}.
\]
The factorials give uniform convergence on every bounded interval
$|x|\le R$ (Weierstrass M-test with the bounds $R^{2n}/n!$ and
$R^{2n+1}/(2n+1)!$), justifying termwise integration when $T$ or $J$ is
applied to the sums. Direct substitutions for (b),(c) were
evaluated explicitly in Exercise~21.2.3.


\Needspace{5\baselineskip}
\par\bigskip
\noindent\textbf{21.3.2}\quad Solve
\[
\psi(x)=x+\int_{0}^{1}(1+xt)\psi(t)\,dt
\]
\erratum{Investigate the following methods:}
\begin{enumerate}
\item[(a)] \erratum{Show that the ordinary Neumann series diverges,
so it cannot supply the solution here.}
\item[(b)] \erratum{Solve by} the separable-kernel technique.
\item[(c)] \erratum{Solve by} educated guessing.
\end{enumerate}
\solution
Let $(Kg)(x)=\int_0^1(1+xt)g(t)dt$.

(a) The ordinary Neumann series prescribed by $\psi=x+K\psi$ is
\[
x+Kx+K^2x+\cdots.
\]
For $g(x)=A+Bx$, since $\int_0^1(A+Bt)dt=A+\frac B2$ and
$\int_0^1t(A+Bt)dt=\frac A2+\frac B3$,
\[
Kg(x)=A+\frac B2+x\left(\frac A2+\frac B3\right),
\]
so (with $A=0$, $B=1$, and then $A=\frac12$, $B=\frac13$)
\[
Kx=\frac12+\frac x3,\qquad
K^2x=\frac12+\frac16+x\left(\frac14+\frac19\right)
=\frac23+\frac{13x}{36}.
\]
Write $K^nx=A_n+B_nx$. Its coefficients obey
\[
A_{n+1}=A_n+\frac{B_n}{2},\qquad
B_{n+1}=\frac{A_n}{2}+\frac{B_n}{3},
\quad (A_0,B_0)=(0,1).
\]
The recurrence has nonnegative coefficients, so by induction all
$A_n,B_n\ge0$. Then $A_{n+1}-A_n=B_n/2\ge0$, so $A_n\ge A_1=1/2$
for every $n\ge1$. Thus $K^nx\ge A_n\ge1/2$ on $[0,1]$ for $n\ge1$:
the terms do not even tend to zero, and the series diverges.

\erratum{The textbook's part (a) asks for the solution by the Neumann
series, but for the equation as printed that series diverges, as just
shown, although the solution exists (parts (b) and (c)). The
Instructor's Manual's answer $-2$ does not address this obstruction.}
Introducing a parameter $z$ locates the obstruction. Since the range of
$K$ consists of linear functions, any solution of $\psi_z=x+zK\psi_z$
has the form $\psi_z=A+Bx$, and by the formula for $Kg$ its
coefficient vector satisfies
\[
\begin{pmatrix}A\\B\end{pmatrix}
=\begin{pmatrix}0\\1\end{pmatrix}
+zM\begin{pmatrix}A\\B\end{pmatrix},\qquad
M=\begin{pmatrix}1&1/2\\1/2&1/3\end{pmatrix}.
\]
Here $D(z)=\det(I-zM)=(1-z)(1-\frac z3)-\frac{z^2}4
=1-\frac43z+\frac1{12}z^2$, and Cramer's rule gives
$A=z/(2D)$, $B=(1-z)/D$:
\[
\psi_z(x)=
\frac{z/2+(1-z)x}{1-\frac43z+\frac1{12}z^2}.
\]
Since $(A_n,B_n)^T=M^n(0,1)^T$, the expansion
$(I-zM)^{-1}=\sum_nz^nM^n$ shows that the Taylor series of $\psi_z$ in
$z$ is the Neumann series $\sum_nz^nK^nx$. The zeros of $D$ are the
roots $z=8\mp2\sqrt{13}$ of $z^2-16z+12=0$ (the characteristic values
of Exercise~21.2.14), and the numerator $z/2$ of $A$ does not vanish
there. So the radius of convergence is $8-2\sqrt{13}\approx0.789<1$,
and the series diverges at $z=1$. Evaluating the rational
continuation at $z=1$, where $D(1)=1-\frac43+\frac1{12}=-\frac14$,
gives $\psi_1=\frac12/(-\frac14)=-2$, but that continuation is not
the sum of the divergent series at $z=1$.

(b) Since $1+xt=1\cdot1+x\cdot t$, the integral term is $c_0+c_1x$
with $c_0=\int_0^1\psi(t)dt$ and $c_1=\int_0^1t\psi(t)dt$.
Separability therefore forces $\psi=A+Bx$ with $A=c_0$, $B=1+c_1$.
By the moments computed in (a), $c_0=A+\frac B2$ and
$c_1=\frac A2+\frac B3$, so comparing coefficients gives
\[
A=A+\frac B2,\qquad
B=1+\frac A2+\frac B3.
\]
The first equation gives $B=0$ and the second then gives $A=-2$.
Thus $\boxed{\psi(x)=-2}$.

(c) Guess a constant $c$. Since $\int_0^1(1+xt)dt=1+\frac x2$,
substitution produces
\[
c=x+c\int_0^1(1+xt)dt=x+c\left(1+\frac x2\right)=c+x(1+c/2).
\]
The coefficient of $x$ vanishes exactly when $c=-2$.
Finally, $x+\int_0^1(1+xt)(-2)dt=x-2-x=-2$,
an explicit check of the solution of the original equation.


\Needspace{5\baselineskip}
\par\bigskip
\noindent\textbf{21.3.3}\quad Solve
\[
\varphi(x)=1+\lambda^{2}\int_{0}^{x}(x-t)\varphi(t)\,dt
\]
by each of the following methods:
\begin{enumerate}
\item[(a)] Reduction to an ODE (find the boundary conditions),
\item[(b)] The Neumann series,
\item[(c)] The use of Laplace transforms.
\end{enumerate}
\par\nobreak\hfill\emph{ANS.}\enspace$\varphi(x)=\cosh\lambda x$.
\solution
(a) Differentiate the equation with the Leibniz rule
$\frac{d}{dx}\int_0^xg(x,t)dt=g(x,x)+\int_0^x\partial_xg(x,t)dt$;
the endpoint term $(x-x)\varphi(x)$ is zero:
\[
\varphi'(x)=\lambda^2\int_0^x\varphi(t)dt,\qquad
\varphi''(x)=\lambda^2\varphi(x).
\]
Setting $x=0$ in the equation and in $\varphi'$ (both integrals
vanish) gives the initial data $\varphi(0)=1$, $\varphi'(0)=0$.
For $\lambda\ne0$, write $\varphi=Ae^{\lambda x}+Be^{-\lambda x}$.
Then $A+B=1$ and $\lambda(A-B)=0$, so $A=B=1/2$.
For $\lambda=0$ the original equation directly gives $\varphi=1$.
In both cases
\[
\boxed{\varphi(x)=\cosh(\lambda x)}.
\]

(b) With $J$ from Exercise~21.3.1, the equation is
$\varphi=1+\lambda^2J\varphi$. The monomial formula of that exercise
with $p=2n$ gives
\[
J\!\left(\frac{t^{2n}}{(2n)!}\right)
=\frac{x^{2n+2}}{(2n)!\,(2n+1)(2n+2)}
=\frac{x^{2n+2}}{(2n+2)!},
\]
so, starting with $J^01=1$, induction gives $J^n1=x^{2n}/(2n)!$, and
the Neumann series is
\[
\sum_{n=0}^\infty\lambda^{2n}J^n1
=\sum_{n=0}^\infty\frac{(\lambda x)^{2n}}{(2n)!}
=\cosh(\lambda x).
\]
It converges uniformly on every finite $x$ interval for every fixed
complex $\lambda$.

(c) Since $\mathcal L\{1\}=1/s$ and $\mathcal L\{x\}=1/s^2$, the
convolution theorem turns $\int_0^x(x-t)\varphi(t)dt$ into
$\Phi(s)/s^2$, so
\[
\Phi(s)=\frac1s+\frac{\lambda^2}{s^2}\Phi(s),\qquad
\Phi(s)=\frac{s}{s^2-\lambda^2}
\]
(multiply by $s^2$: $(s^2-\lambda^2)\Phi=s$).
For $\lambda\ne0$,
\[
\frac{s}{s^2-\lambda^2}
=\frac12\left(\frac1{s-\lambda}+\frac1{s+\lambda}\right),
\]
whose inverse, by $\mathcal L\{e^{ax}\}=1/(s-a)$, is
$\frac12(e^{\lambda x}+e^{-\lambda x})=\cosh(\lambda x)$. At
$\lambda=0$ the transform is $1/s$, the transform of $1=\cosh0$.
As a direct check, for $\lambda\ne0$ integration by parts
(differentiate $x-t$, integrate $\cosh\lambda t$) gives
\[
\begin{aligned}
\int_0^x(x-t)\cosh(\lambda t)dt
&=\left[\frac{(x-t)\sinh\lambda t}{\lambda}\right]_0^x
+\frac1\lambda\int_0^x\sinh\lambda t\,dt\\
&=\frac{\cosh(\lambda x)-1}{\lambda^2},
\end{aligned}
\]
so the right side of the equation is
$1+(\cosh\lambda x-1)=\cosh\lambda x$. As $\lambda\to0$ the integral
tends to $x^2/2=\int_0^x(x-t)dt$, consistent with the zero-parameter
case.


\Needspace{5\baselineskip}
\par\bigskip
\noindent\textbf{21.3.4}
\begin{enumerate}
\item[(a)] In Eq.~(21.59), take $V=V_{0}$, independent of $t$.
Without using Eq.~(21.60), show that Eq.~(21.59) leads directly
to
\[
U(t-t_{0})=\exp\left[-\frac{i}{\hbar}(t-t_{0})V_{0}\right].
\]
\item[(b)] Repeat for Eq.~(21.60) without using Eq.~(21.59).
\end{enumerate}
\erratum{For the power-series derivations, take $V_0$ to be a scalar,
matrix, or bounded operator. State the qualification for an unbounded
self-adjoint operator.}
\solution
(a) With $V(t)=V_0$, Eq.~(21.59) is
\[
U(t,t_0)=I-\frac{i}{\hbar}\int_{t_0}^tV_0U(s,t_0)ds.
\]
Differentiating with respect to $t$ (fundamental theorem of calculus),
and setting $t=t_0$, where the integral vanishes, gives
\[
\partial_tU=-\frac{i}{\hbar}V_0U,\qquad U(t_0,t_0)=I.
\]
Put $\tau=t-t_0$ and $U=\sum_{n\ge0}C_n\tau^n$, so that
$\partial_tU=\sum_{n\ge0}(n+1)C_{n+1}\tau^n$.
Equating powers gives
\[
C_0=I,\qquad (n+1)C_{n+1}
=-\frac{i}{\hbar}V_0C_n.
\]
Induction yields
$C_n=(-iV_0/\hbar)^n/n!$. This series converges for every $\tau$
(its terms are bounded in norm by $(\|V_0\|\,|\tau|/\hbar)^n/n!$) and
solves the initial-value problem, whose solution is unique, so
\[
\boxed{U(t,t_0)=
\exp[-i(t-t_0)V_0/\hbar]}.
\]
It depends only on the time difference because the constant-coefficient
initial-value solution just derived has that form.

(b) Independently, Eq.~(21.60) is the ordered integral series
\[
U=I+\sum_{n=1}^\infty
\left(-\frac{i}{\hbar}\right)^n
\int_{t_0}^t ds_1\int_{t_0}^{s_1}ds_2\cdots
\int_{t_0}^{s_{n-1}}ds_n\,
V(s_1)V(s_2)\cdots V(s_n).
\]
For constant $V_0$ the product is $V_0^n$, which comes out of the
integrals. The remaining scalar simplex integral $S_n(\tau)$, a
function of $\tau=t-t_0$, has $S_0=1$, and in $S_{n+1}$ the inner $n$
integrals (over $s_2,\dots,s_{n+1}$) form $S_n(s_1-t_0)$. With
$u=s_1-t_0$ and the induction hypothesis $S_n(u)=u^n/n!$,
\[
S_{n+1}(\tau)=\int_0^\tau S_n(u)du
=\int_0^\tau\frac{u^n}{n!}du
=\frac{\tau^{n+1}}{(n+1)!}.
\]
Substituting $S_n(\tau)=\tau^n/n!$,
\[
U=I+\sum_{n=1}^\infty\left(-\frac{i}{\hbar}\right)^nV_0^n
\frac{\tau^n}{n!}
=\sum_{n=0}^\infty\frac1{n!}\left(-\frac{i\tau V_0}{\hbar}\right)^n,
\]
precisely the same exponential, without using
Eq.~(21.59). Oriented integrals give the formula also for $\tau<0$.
For a scalar, matrix, or bounded operator $V_0$ both series converge
on finite time intervals.
\erratum{For an unbounded self-adjoint quantum operator, the exponential
is defined by its spectral calculus; an operator power series on every
state cannot be assumed without domain and convergence conditions.}


\Needspace{8\baselineskip}
\subsection*{21.4 Hilbert--Schmidt Theory}

\Needspace{5\baselineskip}
\par\bigskip
\noindent\textbf{21.4.1}\quad In the Fredholm equation
\[
\varphi(x)=\lambda\int_{a}^{b}K(x,t)\varphi(t)\,dt,
\]
assume that the kernel $K(x,t)$ is self-adjoint or Hermitian:
\[
K(x,t)=K^{*}(t,x).
\]
\erratum{Assume a Hilbert--Schmidt kernel,
$\int_a^b\!\int_a^b|K(x,t)|^2\,dx\,dt<\infty$, and nonzero
$L^2$ eigenfunctions.}
Extend the analysis of the present section to show that
\begin{enumerate}
\item[(a)] the eigenfunctions are orthogonal, in the sense that
\[
\int_{a}^{b}\varphi_{m}^{*}(x)\varphi_{n}(x)\,dx=0,
\qquad m\ne n\ (\lambda_{m}\ne\lambda_{n}).
\]
\item[(b)] the eigenvalues are real.
\end{enumerate}
\solution
Let $\langle u,v\rangle=\int_a^b u(x)^*v(x)dx$ and
$(Kg)(x)=\int_a^bK(x,t)g(t)dt$, and assume the
kernel defines a self-adjoint operator on this $L^2$ space (a Hermitian
Hilbert--Schmidt kernel is sufficient). The Hilbert--Schmidt condition
in the statement is what makes the double integral below absolutely
convergent for $L^2$ functions (by the Cauchy--Schwarz inequality), so
that the order of integration may be interchanged (Fubini's theorem).
Interchanging it and using $K(x,t)=K^*(t,x)$ gives the
double-integral identity
\[
\begin{aligned}
\langle\varphi_m,K\varphi_n\rangle
&=\int_a^b\!\!\int_a^b
\varphi_m(x)^*K(x,t)\varphi_n(t)\,dt\,dx\\
&=\int_a^b\left[\int_a^bK(t,x)\varphi_m(x)dx\right]^*
\varphi_n(t)dt
=\langle K\varphi_m,\varphi_n\rangle.
\end{aligned}
\]
A nonzero eigenfunction in the convention
$\varphi_n=\lambda_nK\varphi_n$ has $\lambda_n\ne0$ (otherwise
$\varphi_n=0$), so $K\varphi_n=\varphi_n/\lambda_n$. Inserting this on
both sides of the identity, and taking the factor $1/\lambda_m$ out of
the first (conjugate-linear) slot as $1/\lambda_m^*$,
\[
\frac1{\lambda_n}\langle\varphi_m,\varphi_n\rangle
=\frac1{\lambda_m^*}\langle\varphi_m,\varphi_n\rangle,
\]
or, multiplying by $\lambda_m^*\lambda_n$,
\[
(\lambda_m^*-\lambda_n)
\langle\varphi_m,\varphi_n\rangle=0.
\]
Take $m=n$. Since $\|\varphi_n\|^2>0$, this gives
$\boxed{\lambda_n^*=\lambda_n}$, proving (b).
Now for $\lambda_m\ne\lambda_n$ the factor
$\lambda_m^*-\lambda_n=\lambda_m-\lambda_n$ (by (b)) is nonzero, so
\[
\boxed{\langle\varphi_m,\varphi_n\rangle=0},
\]
which proves (a). Within a degenerate eigenspace orthogonality need
not hold automatically, but Gram--Schmidt produces an orthonormal
basis without changing its eigenvalue.


\Needspace{5\baselineskip}
\par\bigskip
\noindent\textbf{21.4.2}
\begin{enumerate}
\item[(a)] Show that the eigenfunctions of Exercise~21.2.12 are
orthogonal.
\item[(b)] Show that the eigenfunctions of Exercise~21.2.14 are
orthogonal.
\end{enumerate}
\solution
Both kernels are real and symmetric, and in each case the eigenvalues
are distinct, so Exercise~21.4.1 predicts the result; here it is
checked directly.

(a) Use the exact eigenfunctions found in Exercise~21.2.12:
\[
y_0=x,\qquad y_+=x^2+\frac1{\sqrt5},\qquad
y_-=x^2-\frac1{\sqrt5}.
\]
The products $y_0y_\pm=x^3\pm x/\sqrt5$ are odd, so
$\int_{-1}^1y_0y_\pm dx=0$. For the two
even functions, $y_+y_-=x^4-\frac15$ and
\[
\int_{-1}^1y_+y_-dx
=\int_{-1}^1\left(x^4-\frac15\right)dx
=\frac25-\frac25=0.
\]
All three distinct eigenfunctions are therefore mutually orthogonal.

(b) Write $\varphi_\pm=1+b_\pm x$, where
$b_\pm=(-2\pm\sqrt{13})/3$ (Exercise~21.2.14). Their sum and product are
\[
b_++b_-=-\frac43,\qquad b_+b_-=\frac{(-2)^2-13}{9}=-1.
\]
It follows that
\[
\begin{aligned}
\int_0^1\varphi_+\varphi_-dx
&=\int_0^1\bigl[1+(b_++b_-)x+b_+b_-x^2\bigr]dx\\
&=1+\frac{b_++b_-}{2}+\frac{b_+b_-}{3}
=1-\frac23-\frac13=\boxed{0}.
\end{aligned}
\]
Using the exact coefficients is necessary to see exact zero;
rounded decimal eigenfunctions only approximate this cancellation.


\Needspace{5\baselineskip}
\par\bigskip
\noindent\textbf{21.4.3}\quad Use the Hilbert--Schmidt method to
solve the inhomogeneous integral equation
\[
\varphi(x)=x+\frac{1}{2}\int_{-1}^{1}(t+x)\varphi(t)\,dt.
\]
The corresponding homogeneous integral equation was treated in
Example~21.2.3.
\par\emph{Note.}~The application of the Hilbert--Schmidt
technique here is somewhat like using a shotgun to kill a
mosquito, especially when the equation can be solved quickly by
expanding in Legendre polynomials.
\solution
For $K(x,t)=t+x$, choose normalized eigenfunctions
\[
u_\pm(x)=\frac12(1\pm\sqrt3\,x).
\]
Their norms are one, because the odd term of
$(1\pm\sqrt3\,x)^2=1\pm2\sqrt3\,x+3x^2$ integrates to zero and
$\int_{-1}^1(1+3x^2)dx/4=(2+2)/4=1$; their mutual inner product is
$\int_{-1}^1(1-3x^2)dx/4=0$. Evaluate the kernel action; the odd
integrands $t$ and $\pm\sqrt3\,xt$ integrate to zero, leaving
\[
\begin{aligned}
Ku_\pm(x)
&=\int_{-1}^1(t+x)\frac{1\pm\sqrt3\,t}{2}dt
=\frac12\left(\pm\sqrt3\int_{-1}^1t^2dt+x\int_{-1}^1dt\right)\\
&=x\pm\frac1{\sqrt3}
=\pm\frac2{\sqrt3}u_\pm(x).
\end{aligned}
\]
Thus $u_\pm=\lambda_\pm Ku_\pm$ requires $\lambda_\pm(\pm2/\sqrt3)=1$:
the characteristic values are
$\lambda_\pm=\pm\sqrt3/2$, as in Example~21.2.3.

The inhomogeneous term $f(x)=x$ has coefficients
\[
b_\pm=\langle u_\pm,x\rangle
=\pm\frac{\sqrt3}{2}\int_{-1}^1x^2dx
=\pm\frac{\sqrt3}{2}\frac23
=\pm\frac1{\sqrt3}.
\]
Since $x=(u_+-u_-)/\sqrt3$ and every $K\varphi$ is a combination of
$1$ and $x$, the solution of $\varphi=x+\lambda K\varphi$ lies in the
span of $u_\pm$: $\varphi=a_+u_++a_-u_-$. Projecting the equation onto
$u_\pm$ and using $\langle u_\pm,K\varphi\rangle
=\langle Ku_\pm,\varphi\rangle=a_\pm/\lambda_\pm$ ($K$ is real and
symmetric) gives $a_\pm=b_\pm+\lambda a_\pm/\lambda_\pm$. For the
requested coupling $\lambda=1/2$, this Hilbert--Schmidt
projection gives
\[
a_\pm=\frac{b_\pm}{1-\lambda/\lambda_\pm},\qquad
\frac{\lambda}{\lambda_\pm}=\pm\frac1{\sqrt3}.
\]
Specifically,
\[
a_+=\frac{1/\sqrt3}{1-1/\sqrt3}=\frac1{\sqrt3-1}=\frac{\sqrt3+1}{2},\qquad
a_-=\frac{-1/\sqrt3}{1+1/\sqrt3}=-\frac1{\sqrt3+1}=\frac{1-\sqrt3}{2}.
\]
Then $a_++a_-=1$ and $a_+-a_-=\sqrt3$. Reconstructing,
\[
\varphi=a_+u_++a_-u_-
=\frac{a_++a_-}{2}
+\frac{\sqrt3(a_+-a_-)}2x
=\frac12+\frac32x
=\boxed{\frac{1+3x}{2}}.
\]
Directly $\int_{-1}^1\varphi(t)dt=1$ and
$\int_{-1}^1t\varphi(t)dt=1$, so the right side of the
original equation is $x+(1+x)/2=(1+3x)/2$, as required.


\Needspace{5\baselineskip}
\par\bigskip
\noindent\textbf{21.4.4}\quad The Fredholm integral equation
\[
\varphi(x)=\lambda\int_{0}^{\infty}e^{-xt}\varphi(t)\,dt,
\qquad\erratum{x>0},
\]
has an infinite number of solutions, of which one is
\[
\varphi(x)=x^{-1/2},\qquad\lambda=\pi^{-1/2}.
\]
Verify that this is a solution and that it is not normalizable.
\par\emph{Note.}~\erratum{These are improper-integral solutions,
not normalizable Hilbert--Schmidt eigenfunctions. The kernel fails the
Hilbert--Schmidt integrability condition; an infinite interval alone
would not establish that failure. Its Taylor series cannot be integrated
term by term against the displayed power solution.}
\solution
The restriction $x>0$ in the statement is needed because for $x\le0$
the integrand $e^{-xt}t^{-1/2}$ is at least $t^{-1/2}$, whose integral
over $(0,\infty)$ diverges.
For $x>0$, set $u=xt$, so $dt=du/x$ and
$t^{-1/2}=x^{1/2}u^{-1/2}$. Then, with
$\int_0^\infty e^{-u}u^{-1/2}du=\Gamma(\frac12)=\sqrt\pi$,
\[
\int_0^\infty e^{-xt}t^{-1/2}dt
=x^{-1/2}\int_0^\infty e^{-u}u^{-1/2}du
=\sqrt\pi\,x^{-1/2}.
\]
Multiplication by $\lambda=\pi^{-1/2}$ gives
$\boxed{\varphi(x)=x^{-1/2}}$, verifying the equation.

Its truncated squared norm is
\[
\int_\epsilon^R|\varphi(x)|^2dx
=\int_\epsilon^R\frac{dx}{x}
=\ln R-\ln\epsilon.
\]
It diverges as $R\to\infty$ and as $\epsilon\downarrow0$.
No nonzero constant multiple is normalizable on $(0,\infty)$.

\erratum{This is a pointwise improper-integral solution for $x>0$,
not an $L^2$ eigenfunction to which the chapter's compact
Hilbert--Schmidt conclusions automatically apply.}
The obstruction to that particular hypothesis can be checked directly:
since $\int_0^\infty e^{-2xt}dx=1/(2t)$ for $t>0$,
\[
\int_0^\infty\!\!\int_0^\infty|e^{-xt}|^2dx\,dt
=\int_0^\infty\frac{dt}{2t}=\infty,
\]
the last integral diverging logarithmically at both ends.
Thus the kernel is not Hilbert--Schmidt; this calculation identifies
the failed hypothesis more precisely than the infinite interval alone.
Nor can its Taylor series be integrated term by term against $t^{-1/2}$:
the coefficient of $(-x)^m/m!$ would contain
$\int_0^\infty t^{m-1/2}dt$, which diverges for every integer $m\ge0$.

The statement about infinitely many improper-integral solutions can also
be verified explicitly. The same substitution $u=xt$ gives, for $s>0$
and $x>0$, the gamma integral
$\int_0^\infty e^{-xt}t^{s-1}dt=\Gamma(s)x^{-s}$. For $0<a<1$, the
choices $s=1-a$ and $s=a$ give
\[
K(t^{-a})(x)=\Gamma(1-a)x^{a-1},\qquad
K(t^{a-1})(x)=\Gamma(a)x^{-a}.
\]
Put $A=\sqrt{\Gamma(a)}$, $B=\sqrt{\Gamma(1-a)}$, and
$\mu=AB$. Since $\Gamma(a)=A^2$ and $\Gamma(1-a)=B^2$,
\[
K(A t^{-a}\pm B t^{a-1})(x)
=AB^2x^{a-1}\pm BA^2x^{-a}
=\pm\mu(Ax^{-a}\pm Bx^{a-1}).
\]
So $\varphi=\lambda K\varphi$ holds for these combinations when
$\lambda(\pm\mu)=1$. Therefore
\[
\varphi_{a,\pm}(x)=\sqrt{\Gamma(a)}x^{-a}
\pm\sqrt{\Gamma(1-a)}x^{a-1},\qquad
\lambda_{a,\pm}=\pm\frac1\mu=\pm\sqrt{\frac{\sin\pi a}{\pi}},
\]
where the reflection formula $\Gamma(a)\Gamma(1-a)=\pi/\sin\pi a$
gives $\mu^2=\pi/\sin\pi a$.
The minus combination at $a=1/2$ is zero and must be excluded;
the plus combination there is a nonzero multiple of the requested
$x^{-1/2}$. These power combinations belong to the improper-integral
solution class, not to the normalizable $L^2$ eigenfunctions.


\Needspace{5\baselineskip}
\par\bigskip
\noindent\textbf{21.4.5}\quad Given
\[
y(x)=x+\lambda\int_{0}^{1}xt\,y(t)\,dt:
\]
\begin{enumerate}
\item[(a)] Determine $y(x)$ as a Neumann series.
\item[(b)] Find the range of $\lambda$ for which your Neumann
series solution is convergent. Compare with the value obtained
from
\[
|\lambda|\cdot|K|_{\max}<1.
\]
\item[(c)] Find the eigenvalue and the eigenfunction of the
corresponding homogeneous integral equation.
\item[(d)] By the separable-kernel method show that the solution
is
\[
y(x)=\frac{3x}{3-\lambda},\qquad\erratum{\lambda\ne3}.
\]
\item[(e)] Find $y(x)$ by the Hilbert--Schmidt method.
\erratum{Also decide whether a solution exists at $\lambda=3$.}
\end{enumerate}
\solution
Let $(Kg)(x)=x\int_0^1t\,g(t)dt$. Its action on the
inhomogeneous term is $Kx=x\int_0^1t^2dt=x/3$, so $K(x/3^n)=x/3^{n+1}$
and induction gives $K^nx=x/3^n$.

(a) The Neumann series $\sum_n\lambda^nK^nx$ is therefore
\[
\boxed{y(x)=x\sum_{n=0}^\infty(\lambda/3)^n}.
\]
(b) Apart from the identically zero value at $x=0$, this geometric
series converges precisely when $\boxed{|\lambda|<3}$.
For $|\lambda|=3$ its terms $(\lambda/3)^nx$ have the constant modulus
$x$, so they do not tend to zero, and for larger
modulus they grow. On $[0,1]^2$, $|K|_{\max}=\max xt=1$ and the
interval length is one, so the stated sufficient bound only gives
$|\lambda|<1$. It is sufficient, not necessary.

(c) A nonzero homogeneous solution
$y=\lambda x\int_0^1ty(t)dt$ is a multiple of $x$.
Putting $y=Cx$ gives $Cx=\lambda Cx/3$, hence
\[
\boxed{\lambda_1=3,\qquad y_1(x)=Cx,\quad C\ne0}.
\]
(d) For the inhomogeneous equation put
$c=\int_0^1ty(t)dt$, so $y=x(1+\lambda c)$. Taking its moment,
$c=(1+\lambda c)\int_0^1t^2dt$, that is,
\[
c=\frac{1+\lambda c}{3},\qquad
(3-\lambda)c=1.
\]
For $\lambda\ne3$,
\[
\boxed{y(x)=x\left(1+\frac{\lambda}{3-\lambda}\right)
=\frac{3x}{3-\lambda}}.
\]
\erratum{The formula excludes $\lambda=3$. At that value the moment
equation becomes $0=1$, so the inhomogeneous equation has no solution.
The rational solution at other $\lambda$ exists even when its
Neumann series diverges.}

(e) Normalize the homogeneous eigenfunction:
$u(x)=\sqrt3\,x$, since $\int_0^1u^2dx=3\int_0^1x^2dx=1$.
The coefficient of $f=x$ is
$b=\int_0^1u(x)x\,dx=\sqrt3/3=1/\sqrt3$.
Since $f=u/\sqrt3$ lies in the span of $u$, and so does every $Ky$,
the solution is $y=au$. Projecting $y=x+\lambda Ky$ onto $u$ and
using $\langle u,Ky\rangle=\langle Ku,y\rangle=a/\lambda_1$ gives
$a=b+\lambda a/\lambda_1$, i.e.\ $(\lambda_1-\lambda)a=b\lambda_1$.
For $\lambda\ne\lambda_1=3$, Hilbert--Schmidt projection thus gives
\[
y(x)=\frac{b\lambda_1}{\lambda_1-\lambda}u(x)
=\frac{(1/\sqrt3)3}{3-\lambda}\sqrt3\,x
=\frac{3x}{3-\lambda}.
\]
At $\lambda=\lambda_1=3$ the projected equation reads
$0\cdot a=3b=\sqrt3\ne0$, because the overlap $b$ is nonzero: there is
no solution, in agreement with (d).


\Needspace{5\baselineskip}
\par\bigskip
\noindent\textbf{21.4.6}\quad In Exercise~21.2.11 it was found
that the integral equation
\[
\varphi(x)=\lambda\int_{0}^{2\pi}\cos(x-t)\varphi(t)\,dt
\]
had (unnormalized) eigenfunctions $\cos x$ and $\sin x$, both with
eigenvalue $\lambda_{i}=1/\pi$. Show that the kernel of this
integral equation has an expansion of the form
\[
K(x,t)=\sum_{n=1}^{2}\frac{\varphi_{n}(x)\varphi_{n}(t)}
{\lambda_{n}},
\]
\erratum{where $\varphi_n$ in this expansion are the unit-norm
versions of the stated eigenfunctions.}
\solution
The squared norms of the stated functions are (Exercise~21.2.11;
$\cos^2x,\sin^2x=(1\pm\cos2x)/2$, and $\cos2x$ integrates to zero
over a period)
\[
\int_0^{2\pi}\cos^2x\,dx
=\int_0^{2\pi}\sin^2x\,dx=\pi,
\]
and their mutual integral is
$\int_0^{2\pi}\sin x\cos x\,dx=\frac12\int_0^{2\pi}\sin2x\,dx=0$.
Therefore use
\[
u_1(x)=\frac{\cos x}{\sqrt\pi},\qquad
u_2(x)=\frac{\sin x}{\sqrt\pi}.
\]
With $\lambda_1=\lambda_2=1/\pi$ and the addition theorem for the
cosine,
\[
\begin{aligned}
\sum_{n=1}^2\frac{u_n(x)u_n(t)}{\lambda_n}
&=\pi\left[\frac{\cos x\cos t}{\pi}
+\frac{\sin x\sin t}{\pi}\right]\\
&=\cos x\cos t+\sin x\sin t
=\boxed{\cos(x-t)=K(x,t)}.
\end{aligned}
\]
\erratum{The functions used in the displayed kernel expansion must
be normalized. Substituting the initially stated unnormalized
$\cos x,\sin x$ would instead give $\pi K(x,t)$.}
Equivalently, retain unnormalized functions but divide each term by
its squared norm as well as by its characteristic value.


\Needspace{5\baselineskip}
\par\bigskip
\noindent\textbf{21.4.7}\quad The integral equation
\[
\varphi(x)=\lambda\int_{0}^{1}(1+xt)\varphi(t)\,dt
\]
has eigenvalues \erratum{$\lambda_{1,2}=8\mp2\sqrt{13}$}
(approximately $0.7889$ and $15.211$). The corresponding
unnormalized eigenfunctions are
\erratum{$\varphi_{1,2}(x)=1+(-2\pm\sqrt{13})x/3$},
approximately $1+0.5352x$ and $1-1.8685x$.
\begin{enumerate}
\item[(a)] Show that these eigenfunctions are orthogonal over the
interval $[0,1]$.
\item[(b)] Normalize the eigenfunctions to unity.
\item[(c)] \erratum{Using the normalized functions from part (b),} show that
\[
K(x,t)=\frac{\varphi_{1}(x)\varphi_{1}(t)}{\lambda_{1}}
+\frac{\varphi_{2}(x)\varphi_{2}(t)}{\lambda_{2}}.
\]
\end{enumerate}
\par\nobreak\hfill\emph{ANS.}\enspace
(b) $\varphi_{1}(x)=0.7831+0.4191x$,\par\nobreak\hfill
$\varphi_{2}(x)=1.8403-3.4386x$.
\solution
Use the exact values behind the displayed decimals:
\[
b_\pm=\frac{-2\pm\sqrt{13}}3,\qquad
\lambda_\pm=8\mp2\sqrt{13},\qquad
v_\pm(x)=1+b_\pm x.
\]
\erratum{The textbook's decimal eigenpairs are approximations, so
exact orthogonality and exact kernel reconstruction use these radical
values. In part (c) the eigenfunctions must be the unit-norm ones of
part (b): with the unnormalized $v_\pm$ the same sum is
$\frac43+\frac59(x+t)+\frac{16}{27}xt$, not $K(x,t)$.}

(a) As in Exercise~21.4.2(b), since $b_++b_-=-4/3$ and $b_+b_-=-1$,
\[
\int_0^1v_+(x)v_-(x)dx
=\int_0^1\bigl[1+(b_++b_-)x+b_+b_-x^2\bigr]dx
=1+\frac{b_++b_-}{2}+\frac{b_+b_-}{3}=0.
\]
(b) Their squared norms are
\[
N_\pm^2=\int_0^1(1+2b_\pm x+b_\pm^2x^2)dx
=1+b_\pm+\frac{b_\pm^2}{3}.
\]
Since $3b_\pm+2=\pm\sqrt{13}$, squaring gives the root equation
$9b_\pm^2+12b_\pm-9=0$, i.e.\ $3b_\pm^2+4b_\pm-3=0$, so
$b_\pm^2=1-\frac43b_\pm$. This reduces $N_\pm^2$ to
\[
N_\pm^2=1+b_\pm+\frac13-\frac{4b_\pm}9
=\frac43+\frac{5b_\pm}{9}
=\frac{36+5(-2\pm\sqrt{13})}{27}
=\frac{26\pm5\sqrt{13}}{27}.
\]
Thus the normalized functions are
\[
\boxed{u_\pm(x)=
\sqrt{\frac{27}{26\pm5\sqrt{13}}}
\left(1+\frac{-2\pm\sqrt{13}}3x\right)}.
\]
Their numerical forms ($u_+$ and $u_-$ are the answer's $\varphi_1$
and $\varphi_2$) are
\[
u_+(x)=0.7831024897+0.4191037337x,\qquad
u_-(x)=1.8403126068-3.4386555600x.
\]
\erratum{The printed coefficient $3.4386$ in the answer is not
correctly rounded: $3.43865556\ldots$ rounds to $3.4387$.}

(c) Put $A_\pm=1/(\lambda_\pm N_\pm^2)$. Multiplying out,
\[
\lambda_\pm N_\pm^2
=\frac{(8\mp2\sqrt{13})(26\pm5\sqrt{13})}{27}
=\frac{78\mp12\sqrt{13}}{27}
=\frac{26\mp4\sqrt{13}}{9},
\]
and rationalizing, with $26^2-16\cdot13=468=9\cdot52$,
\[
A_\pm=\frac{9(26\pm4\sqrt{13})}{468}
=\frac{26\pm4\sqrt{13}}{52}
=\frac12\pm\frac1{\sqrt{13}}.
\]
Each kernel term is
$u_\pm(x)u_\pm(t)/\lambda_\pm
=A_\pm(1+b_\pm x)(1+b_\pm t)
=A_\pm\bigl[1+b_\pm(x+t)+b_\pm^2xt\bigr]$,
so in their sum the constant, $x+t$, and $xt$
coefficients are respectively
$\sum A_\pm$, $\sum A_\pm b_\pm$, and
$\sum A_\pm b_\pm^2$. Evaluate each, using
$b_+-b_-=2\sqrt{13}/3$:
\[
\sum A_\pm=1,\qquad
\sum A_\pm b_\pm
=\frac12(b_++b_-)+\frac{b_+-b_-}{\sqrt{13}}
=\frac12\left(-\frac43\right)
+\frac1{\sqrt{13}}\frac{2\sqrt{13}}3=0,
\]
and, using $b_\pm^2=1-\tfrac43b_\pm$,
\[
\sum A_\pm b_\pm^2=1-\frac43(0)=1.
\]
Consequently
\[
\boxed{\frac{u_+(x)u_+(t)}{\lambda_+}
+\frac{u_-(x)u_-(t)}{\lambda_-}
=1+xt}.
\]


\Needspace{5\baselineskip}
\par\bigskip
\noindent\textbf{21.4.8}\quad An alternate form of the solution to
the inhomogeneous integral equation, Eq.~(21.81), is
\[
\varphi(x)=\erratum{f_0(x)+}\sum_{i=1}^{\infty}\frac{b_{i}\lambda_{i}}
{\lambda_{i}-\lambda}\,\varphi_{i}(x).
\]
\erratum{Here $\varphi_i$ are orthonormal eigenfunctions for the
nonzero operator eigenvalues, $b_i=\langle\varphi_i,f\rangle$,
$f_0$ is the projection of $f$ onto $\ker K$, and
$\lambda\ne\lambda_i$. The textbook's pure sum applies when $f_0=0$.}
\begin{enumerate}
\item[(a)] Derive this form without using Eq.~(21.89).
\item[(b)] Show that this form and Eq.~(21.89) are equivalent.
\end{enumerate}
\solution
Let $K$ be the self-adjoint compact integral operator in the
chapter, and let its normalized nonzero-eigenvalue functions satisfy
$K\varphi_i=\varphi_i/\lambda_i$.

(a) By the Hilbert--Schmidt (spectral) theorem for a compact
self-adjoint operator, $L^2$ is the orthogonal sum of $\ker K$ and the
closed span of the $\varphi_i$. Decompose both functions accordingly
into their nullspace components and
their components along the nonzero-eigenvalue eigenfunctions:
\[
f=f_0+\sum_i b_i\varphi_i,\qquad
\varphi=\varphi_0+\sum_i a_i\varphi_i.
\]
The projection of $K\varphi$ onto $\ker K$ is zero, since
$K$ is self-adjoint: for $h\in\ker K$,
$\langle h,K\varphi\rangle=\langle Kh,\varphi\rangle=0$.
Projecting the integral equation onto that nullspace gives
$\varphi_0=f_0$.
Project $\varphi=f+\lambda K\varphi$ onto $\varphi_i$, using
$\langle\varphi_i,K\varphi\rangle=\langle K\varphi_i,\varphi\rangle
=a_i/\lambda_i$ ($\lambda_i$ is real by Exercise~21.4.1):
\[
a_i=b_i+\frac{\lambda}{\lambda_i}a_i,\qquad
(\lambda_i-\lambda)a_i=b_i\lambda_i.
\]
For $\lambda\ne\lambda_i$ this gives
\[
\boxed{\varphi(x)=f_0(x)+\sum_i
\frac{b_i\lambda_i}{\lambda_i-\lambda}\varphi_i(x)},
\]
without using Eq.~(21.89).

(b) Split the scalar factor:
\[
\frac{\lambda_i}{\lambda_i-\lambda}
=1+\frac{\lambda}{\lambda_i-\lambda}.
\]
Summing the first terms gives $\sum_ib_i\varphi_i=f-f_0$; adding the
explicit nullspace term restores $f$, so
\[
\varphi(x)=f(x)+\lambda\sum_i
\frac{\varphi_i(x)}{\lambda_i-\lambda}
\underbrace{\int_a^b\varphi_i(t)^*f(t)dt}_{b_i}.
\]
This is Eq.~(21.89), with the conjugate included for a complex
orthonormal basis.

\erratum{The textbook's pure eigenfunction sum assumes that $f$ has no
component in the nullspace of $K$. Nonzero-eigenvalue eigenfunctions
need not span the whole $L^2$ space, for example for the finite-rank
kernels in this chapter. In general an additional unchanged nullspace
component is required.}
Thus Eq.~(21.89) already retains $f_0$ in its explicit $f$ term,
while the textbook's pure sum requires the extra condition
$f_0=0$.

For a concrete check, take $K(x,t)=xt$ on $[0,1]$,
$f=1$, and $\lambda=1$. Here $u=\sqrt3x$ satisfies
$Ku=x\int_0^1\sqrt3\,t^2dt=u/3$, so $\lambda_1=3$;
$b=\int_0^1\sqrt3\,x\,dx=\sqrt3/2$, and
$f_0=f-bu=1-3x/2$. The corrected formula gives, with
$\frac{3b}{3-1}u=\frac{3\sqrt3}{4}\sqrt3\,x=\frac94x$,
\[
\varphi=f_0+\frac{3b}{3-1}u
=1-\frac32x+\frac94x
=\boxed{1+\frac34x}.
\]
Indeed $1+x\int_0^1t(1+3t/4)dt=1+x(\frac12+\frac14)=1+3x/4$.
The sum without $f_0$ would give $9x/4$, which fails the equation:
its right side is again $1+3x/4$.

For a compact self-adjoint operator the nonzero operator eigenvalues
$\mu_i=1/\lambda_i$ can accumulate only at zero. For a fixed finite
$\lambda$, all but finitely many satisfy $|\lambda\mu_i|<1/2$,
so their multipliers
$\lambda_i/(\lambda_i-\lambda)=1/(1-\lambda\mu_i)$ obey
$|1/(1-\lambda\mu_i)|\le1/(1-|\lambda\mu_i|)\le2$.
The finitely many remaining denominators are nonzero away from
resonance and therefore have a positive minimum modulus. Hence all
multipliers are bounded by some $M$, so $|a_i|\le M|b_i|$, and
Parseval's identity in the closed span of the $\varphi_i$ gives
$\sum_i|a_i|^2\le M^2\sum_i|b_i|^2=M^2\|f-f_0\|^2<\infty$: the
solution expansion converges in $L^2$ whenever $f\in L^2$.
At $\lambda=\lambda_j$, the projected equation instead requires
$b_j=0$ for every eigenfunction in that eigenspace; if so, its
solution coefficients are free, and otherwise there is no solution.
